% Options for packages loaded elsewhere
% Options for packages loaded elsewhere
\PassOptionsToPackage{unicode}{hyperref}
\PassOptionsToPackage{hyphens}{url}
\PassOptionsToPackage{dvipsnames,svgnames,x11names}{xcolor}
%
\documentclass[
  spanish,
  letterpaper,
]{scrbook}
\usepackage{xcolor}
\usepackage[top=25mm,bottom=25mm,left=25mm,right=25mm]{geometry}
\usepackage{amsmath,amssymb}
\setcounter{secnumdepth}{2}
\usepackage{iftex}
\ifPDFTeX
  \usepackage[T1]{fontenc}
  \usepackage[utf8]{inputenc}
  \usepackage{textcomp} % provide euro and other symbols
\else % if luatex or xetex
  \usepackage{unicode-math} % this also loads fontspec
  \defaultfontfeatures{Scale=MatchLowercase}
  \defaultfontfeatures[\rmfamily]{Ligatures=TeX,Scale=1}
\fi
\usepackage{lmodern}
\ifPDFTeX\else
  % xetex/luatex font selection
\fi
% Use upquote if available, for straight quotes in verbatim environments
\IfFileExists{upquote.sty}{\usepackage{upquote}}{}
\IfFileExists{microtype.sty}{% use microtype if available
  \usepackage[]{microtype}
  \UseMicrotypeSet[protrusion]{basicmath} % disable protrusion for tt fonts
}{}
\makeatletter
\@ifundefined{KOMAClassName}{% if non-KOMA class
  \IfFileExists{parskip.sty}{%
    \usepackage{parskip}
  }{% else
    \setlength{\parindent}{0pt}
    \setlength{\parskip}{6pt plus 2pt minus 1pt}}
}{% if KOMA class
  \KOMAoptions{parskip=half}}
\makeatother
% Make \paragraph and \subparagraph free-standing
\makeatletter
\ifx\paragraph\undefined\else
  \let\oldparagraph\paragraph
  \renewcommand{\paragraph}{
    \@ifstar
      \xxxParagraphStar
      \xxxParagraphNoStar
  }
  \newcommand{\xxxParagraphStar}[1]{\oldparagraph*{#1}\mbox{}}
  \newcommand{\xxxParagraphNoStar}[1]{\oldparagraph{#1}\mbox{}}
\fi
\ifx\subparagraph\undefined\else
  \let\oldsubparagraph\subparagraph
  \renewcommand{\subparagraph}{
    \@ifstar
      \xxxSubParagraphStar
      \xxxSubParagraphNoStar
  }
  \newcommand{\xxxSubParagraphStar}[1]{\oldsubparagraph*{#1}\mbox{}}
  \newcommand{\xxxSubParagraphNoStar}[1]{\oldsubparagraph{#1}\mbox{}}
\fi
\makeatother


\usepackage{longtable,booktabs,array}
\usepackage{calc} % for calculating minipage widths
% Correct order of tables after \paragraph or \subparagraph
\usepackage{etoolbox}
\makeatletter
\patchcmd\longtable{\par}{\if@noskipsec\mbox{}\fi\par}{}{}
\makeatother
% Allow footnotes in longtable head/foot
\IfFileExists{footnotehyper.sty}{\usepackage{footnotehyper}}{\usepackage{footnote}}
\makesavenoteenv{longtable}
\usepackage{graphicx}
\makeatletter
\newsavebox\pandoc@box
\newcommand*\pandocbounded[1]{% scales image to fit in text height/width
  \sbox\pandoc@box{#1}%
  \Gscale@div\@tempa{\textheight}{\dimexpr\ht\pandoc@box+\dp\pandoc@box\relax}%
  \Gscale@div\@tempb{\linewidth}{\wd\pandoc@box}%
  \ifdim\@tempb\p@<\@tempa\p@\let\@tempa\@tempb\fi% select the smaller of both
  \ifdim\@tempa\p@<\p@\scalebox{\@tempa}{\usebox\pandoc@box}%
  \else\usebox{\pandoc@box}%
  \fi%
}
% Set default figure placement to htbp
\def\fps@figure{htbp}
\makeatother



\ifLuaTeX
\usepackage[bidi=basic]{babel}
\else
\usepackage[bidi=default]{babel}
\fi
% get rid of language-specific shorthands (see #6817):
\let\LanguageShortHands\languageshorthands
\def\languageshorthands#1{}


\setlength{\emergencystretch}{3em} % prevent overfull lines

\providecommand{\tightlist}{%
  \setlength{\itemsep}{0pt}\setlength{\parskip}{0pt}}



 


\usepackage{emoji}
\makeatletter
\@ifpackageloaded{tcolorbox}{}{\usepackage[skins,breakable]{tcolorbox}}
\@ifpackageloaded{fontawesome5}{}{\usepackage{fontawesome5}}
\definecolor{quarto-callout-color}{HTML}{909090}
\definecolor{quarto-callout-note-color}{HTML}{0758E5}
\definecolor{quarto-callout-important-color}{HTML}{CC1914}
\definecolor{quarto-callout-warning-color}{HTML}{EB9113}
\definecolor{quarto-callout-tip-color}{HTML}{00A047}
\definecolor{quarto-callout-caution-color}{HTML}{FC5300}
\definecolor{quarto-callout-color-frame}{HTML}{acacac}
\definecolor{quarto-callout-note-color-frame}{HTML}{4582ec}
\definecolor{quarto-callout-important-color-frame}{HTML}{d9534f}
\definecolor{quarto-callout-warning-color-frame}{HTML}{f0ad4e}
\definecolor{quarto-callout-tip-color-frame}{HTML}{02b875}
\definecolor{quarto-callout-caution-color-frame}{HTML}{fd7e14}
\makeatother
\makeatletter
\@ifpackageloaded{bookmark}{}{\usepackage{bookmark}}
\makeatother
\makeatletter
\@ifpackageloaded{caption}{}{\usepackage{caption}}
\AtBeginDocument{%
\ifdefined\contentsname
  \renewcommand*\contentsname{Tabla de contenidos}
\else
  \newcommand\contentsname{Tabla de contenidos}
\fi
\ifdefined\listfigurename
  \renewcommand*\listfigurename{Listado de Figuras}
\else
  \newcommand\listfigurename{Listado de Figuras}
\fi
\ifdefined\listtablename
  \renewcommand*\listtablename{Listado de Tablas}
\else
  \newcommand\listtablename{Listado de Tablas}
\fi
\ifdefined\figurename
  \renewcommand*\figurename{Figura}
\else
  \newcommand\figurename{Figura}
\fi
\ifdefined\tablename
  \renewcommand*\tablename{Tabla}
\else
  \newcommand\tablename{Tabla}
\fi
}
\@ifpackageloaded{float}{}{\usepackage{float}}
\floatstyle{ruled}
\@ifundefined{c@chapter}{\newfloat{codelisting}{h}{lop}}{\newfloat{codelisting}{h}{lop}[chapter]}
\floatname{codelisting}{Listado}
\newcommand*\listoflistings{\listof{codelisting}{Listado de Listados}}
\makeatother
\makeatletter
\makeatother
\makeatletter
\@ifpackageloaded{caption}{}{\usepackage{caption}}
\@ifpackageloaded{subcaption}{}{\usepackage{subcaption}}
\makeatother
\usepackage{bookmark}
\IfFileExists{xurl.sty}{\usepackage{xurl}}{} % add URL line breaks if available
\urlstyle{same}
\hypersetup{
  pdftitle={El libro de notas de inglés},
  pdfauthor={Maicel Maicel},
  pdflang={es},
  colorlinks=true,
  linkcolor={Maroon},
  filecolor={Maroon},
  citecolor={Blue},
  urlcolor={Blue},
  pdfcreator={LaTeX via pandoc}}


\title{El libro de notas de inglés}
\usepackage{etoolbox}
\makeatletter
\providecommand{\subtitle}[1]{% add subtitle to \maketitle
  \apptocmd{\@title}{\par {\large #1 \par}}{}{}
}
\makeatother
\subtitle{De A1 a C1 --- El libro que ojalá hubiera existido antes}
\author{Maicel Maicel}
\date{2026-06-01}
\begin{document}
\frontmatter
\maketitle

\renewcommand*\contentsname{Tabla de contenidos}
{
\hypersetup{linkcolor=}
\setcounter{tocdepth}{2}
\tableofcontents
}

\mainmatter
\bookmarksetup{startatroot}

\chapter*{Preface}\label{preface}
\addcontentsline{toc}{chapter}{Preface}

\markboth{Preface}{Preface}

This is a Quarto book.

\emph{Por un hispanohablante, para hispanohablantes}

\bookmarksetup{startatroot}

\chapter{Antes de empezar}\label{antes-de-empezar}

\begin{quote}
\subsection{💬 Antes de empezar}\label{antes-de-empezar-1}

Este libro no es un libro de texto frío y aburrido. Es tu mejor amigo
que ya aprendió inglés y te está contando \textbf{todos los secretos.}

Léelo como si fuera una conversación, no como un manual de
instrucciones. Cada página tiene un propósito: \textbf{que aprendas el
doble en la mitad del tiempo.}
\end{quote}

\begin{center}\rule{0.5\linewidth}{0.5pt}\end{center}

\begin{quote}
\subsection{🗺️ Cómo usar este libro}\label{cuxf3mo-usar-este-libro}

\begin{longtable}[]{@{}ll@{}}
\toprule\noalign{}
Símbolo & Significa \\
\midrule\noalign{}
\endhead
\bottomrule\noalign{}
\endlastfoot
🎯 & La idea central --- léela primero siempre \\
📐 & La fórmula o estructura gramatical \\
💡 & Truco mnemotécnico --- esto lo recordarás siempre \\
✅ & Ejemplos correctos \\
⚠️ & Error típico del hispanohablante \\
🔥 & Tip directo para el examen \\
🧠 & Momento de reflexión profunda \\
📝 & Mini ejercicio mental \\
\end{longtable}
\end{quote}

\begin{center}\rule{0.5\linewidth}{0.5pt}\end{center}

\begin{quote}
\subsection{⚡ Progreso del libro}\label{progreso-del-libro}

\begin{longtable}[]{@{}lll@{}}
\toprule\noalign{}
Nivel & Capítulos & Estado \\
\midrule\noalign{}
\endhead
\bottomrule\noalign{}
\endlastfoot
📘 A1 & 1 --- 4 & Bases absolutas \\
📗 A2 & 5 --- 7 & Supervivencia \\
📙 B1 & 8 --- 13 & Fluidez inicial \\
📕 B2 & 14 --- 22 & Dominio real \\
\end{longtable}
\end{quote}

\part{📘 Parte 1: Nivel A1}

\chapter{El Verbo To Be}\label{el-verbo-to-be}

La arquitectura fundamental del inglés

\hfill\break

El verbo \emph{To Be} es la estructura más frecuente, versátil y
sintácticamente única del inglés. A diferencia de la mayoría de los
verbos ingleses, funciona como su propio operador: no requiere
auxiliares externos (\emph{do} / \emph{does}) para negar o preguntar.

\begin{center}\rule{0.5\linewidth}{0.5pt}\end{center}

\begin{tcolorbox}[enhanced jigsaw, toprule=.15mm, leftrule=.75mm, opacityback=0, coltitle=black, colbacktitle=quarto-callout-note-color!10!white, left=2mm, opacitybacktitle=0.6, breakable, title={🎯 La idea central}, toptitle=1mm, titlerule=0mm, arc=.35mm, colback=white, bottomtitle=1mm, colframe=quarto-callout-note-color-frame, rightrule=.15mm, bottomrule=.15mm]

\emph{To Be} funciona exactamente como un \textbf{signo de igualdad
(\(=\))} o puente copulativo. Su único objetivo es conectar un
\textbf{Sujeto} con un \textbf{Atributo} (una identidad, un estado, una
ubicación o una propiedad).

Mientras el español fragmenta este concepto en dos verbos distintos
(\emph{Ser} y \emph{Estar}), la mente angloparlante procesa ambos
escenarios bajo una única categoría de existencia.

\end{tcolorbox}

\begin{center}\rule{0.5\linewidth}{0.5pt}\end{center}

\section{📐 La fórmula o estructura
gramatical}\label{la-fuxf3rmula-o-estructura-gramatical}

La morfología de \emph{To Be} en presente se divide en tres formas
estructurales fijas según el tipo de sujeto:

\[\text{Sujeto} + \text{Verbo TO BE} + \text{Atributo}\]

\begin{longtable}[]{@{}
  >{\raggedright\arraybackslash}p{(\linewidth - 6\tabcolsep) * \real{0.2222}}
  >{\raggedright\arraybackslash}p{(\linewidth - 6\tabcolsep) * \real{0.2222}}
  >{\centering\arraybackslash}p{(\linewidth - 6\tabcolsep) * \real{0.2778}}
  >{\centering\arraybackslash}p{(\linewidth - 6\tabcolsep) * \real{0.2778}}@{}}
\toprule\noalign{}
\begin{minipage}[b]{\linewidth}\raggedright
Tipo de Sujeto
\end{minipage} & \begin{minipage}[b]{\linewidth}\raggedright
Pronombres
\end{minipage} & \begin{minipage}[b]{\linewidth}\centering
Forma
\end{minipage} & \begin{minipage}[b]{\linewidth}\centering
Contracción Estándar
\end{minipage} \\
\midrule\noalign{}
\endhead
\bottomrule\noalign{}
\endlastfoot
\textbf{Primera persona (Singular)} & I & \textbf{am} & I'm \\
\textbf{Tercera persona (Singular)} & He, She, It & \textbf{is} & He's /
She's / It's \\
\textbf{Plurales y Segunda persona} & You, We, They & \textbf{are} &
You're / We're / They're \\
\end{longtable}

\subsection{El Mecanismo de Inversión
Sintáctica}\label{el-mecanismo-de-inversiuxf3n-sintuxe1ctica}

Para estructurar preguntas, el inglés no depende de la entonación como
el español; exige una \textbf{alerta estructural}. El verbo operador
ejecuta un movimiento sintáctico saltando a la izquierda del sujeto:

\begin{verbatim}

Afirmación:   [ You ]  [ are ]  ready.
│         │
└────┐┌───┘  (Inversión)
vv
Pregunta:     [ Are ]  [ you ]  ready?
\end{verbatim}

\begin{center}\rule{0.5\linewidth}{0.5pt}\end{center}

\begin{tcolorbox}[enhanced jigsaw, toprule=.15mm, leftrule=.75mm, opacityback=0, coltitle=black, colbacktitle=quarto-callout-tip-color!10!white, left=2mm, opacitybacktitle=0.6, breakable, title={💡 Truco mnemotécnico}, toptitle=1mm, titlerule=0mm, arc=.35mm, colback=white, bottomtitle=1mm, colframe=quarto-callout-tip-color-frame, rightrule=.15mm, bottomrule=.15mm]

Para recordar la sintaxis interrogativa, imagina que el verbo \emph{To
Be} actúa como un \textbf{``escudo delantero''}. En una pregunta, el
verbo siempre debe ir al frente de la oración para recibir el impacto
del signo de interrogación al final. Si el sujeto va primero, el cerebro
nativo asume automáticamente que es una declaración, no una pregunta.

\end{tcolorbox}

\begin{center}\rule{0.5\linewidth}{0.5pt}\end{center}

\section{✅ Ejemplos correctos}\label{ejemplos-correctos}

\begin{itemize}
\tightlist
\item
  \textbf{Identidad:} I \textbf{am} a doctor. \(\rightarrow\) \emph{Soy
  médico.}
\item
  \textbf{Ubicación:} She \textbf{is} in London. \(\rightarrow\)
  \emph{Ella está en Londres.}
\item
  \textbf{Estado:} We \textbf{are} ready. \(\rightarrow\) \emph{Estamos
  listos.}
\item
  \textbf{Pregunta con inversión:} \textbf{Are} they late?
  \(\rightarrow\) \emph{¿Llegan tarde?}
\end{itemize}

\begin{center}\rule{0.5\linewidth}{0.5pt}\end{center}

\begin{tcolorbox}[enhanced jigsaw, toprule=.15mm, leftrule=.75mm, opacityback=0, coltitle=black, colbacktitle=quarto-callout-warning-color!10!white, left=2mm, opacitybacktitle=0.6, breakable, title={⚠️ Error típico del hispanohablante}, toptitle=1mm, titlerule=0mm, arc=.35mm, colback=white, bottomtitle=1mm, colframe=quarto-callout-warning-color-frame, rightrule=.15mm, bottomrule=.15mm]

El cerebro bilingüe tiende a transferir los vicios del idioma materno.
Detecta y bloquea estos tres reflejos automáticos:

\begin{enumerate}
\def\labelenumi{\arabic{enumi}.}
\tightlist
\item
  \textbf{El Sujeto Fantasma:} En español existe el sujeto elíptico
  (\emph{``Es importante''}). En inglés, omitir el sujeto destruye la
  sintaxis. Si no hay un agente real, debes usar el pronombre neutro
  \textbf{It}.

  \begin{itemize}
  \tightlist
  \item
    ❌ \emph{Is important.} \(\rightarrow\) ✅ \textbf{It is} important.
  \end{itemize}
\item
  \textbf{La Posesión de la Edad:} En inglés la edad no se posee
  (\emph{Tener}), se \emph{es} viejo a través del tiempo.

  \begin{itemize}
  \tightlist
  \item
    ❌ \emph{I have 30 years.} \(\rightarrow\) ✅ \textbf{I am} 30 years
    old.
  \end{itemize}
\item
  \textbf{Clima Impersonal:}

  \begin{itemize}
  \tightlist
  \item
    ❌ \emph{Is raining.} \(\rightarrow\) ✅ \textbf{It is} raining.
  \end{itemize}
\end{enumerate}

\end{tcolorbox}

\begin{center}\rule{0.5\linewidth}{0.5pt}\end{center}

\begin{tcolorbox}[enhanced jigsaw, toprule=.15mm, leftrule=.75mm, opacityback=0, coltitle=black, colbacktitle=quarto-callout-important-color!10!white, left=2mm, opacitybacktitle=0.6, breakable, title={🔥 Tip directo para el examen}, toptitle=1mm, titlerule=0mm, arc=.35mm, colback=white, bottomtitle=1mm, colframe=quarto-callout-important-color-frame, rightrule=.15mm, bottomrule=.15mm]

Si te enfrentas a una prueba de redacción formal (\textbf{Writing B2/C1}
como Cambridge, IELTS o TOEFL), \textbf{las contracciones están
estrictamente prohibidas}. El uso de \emph{don't}, \emph{isn't} o
\emph{they're} resta autoridad académica a tu texto. * ❌ \emph{The
results \textbf{aren't} conclusive.} (Registro informal/hablado). * ✅
\emph{The results \textbf{are not} conclusive.} (Registro
académico/formal).

\end{tcolorbox}

\begin{center}\rule{0.5\linewidth}{0.5pt}\end{center}

\begin{tcolorbox}[enhanced jigsaw, toprule=.15mm, leftrule=.75mm, opacityback=0, coltitle=black, colbacktitle=quarto-callout-caution-color!10!white, left=2mm, opacitybacktitle=0.6, breakable, title={🧠 Momento de reflexión profunda}, toptitle=1mm, titlerule=0mm, arc=.35mm, colback=white, bottomtitle=1mm, colframe=quarto-callout-caution-color-frame, rightrule=.15mm, bottomrule=.15mm]

¿Por qué el inglés es tan rígido con la presencia del sujeto
(\texttt{It}) mientras que el español lo descarta fácilmente? La
respuesta está en la morfología del verbo.

En español, la terminación del verbo te dice quién habla
(\emph{``Estamo\textbf{s}''} solo puede ser \emph{Nosotros}). En inglés,
formas como \texttt{are} se repiten para \emph{You}, \emph{We} y
\emph{They}. Sin el pronombre explícito delante, el cerebro nativo es
incapaz de descodificar quién ejecuta la acción. El pronombre en inglés
no es un adorno; es un requisito de claridad matemática.

\end{tcolorbox}

\begin{center}\rule{0.5\linewidth}{0.5pt}\end{center}

\begin{tcolorbox}[enhanced jigsaw, toprule=.15mm, leftrule=.75mm, opacityback=0, coltitle=black, colbacktitle=quarto-callout-note-color!10!white, left=2mm, opacitybacktitle=0.6, breakable, title={📝 Mini ejercicio mental}, toptitle=1mm, titlerule=0mm, arc=.35mm, colback=white, bottomtitle=1mm, colframe=quarto-callout-note-color-frame, rightrule=.15mm, bottomrule=.15mm]

\emph{Fuerza a tu mente a generar la solución y justificar su estructura
antes de desplegar la respuesta.}

\begin{enumerate}
\def\labelenumi{\arabic{enumi}.}
\tightlist
\item
  Corrige la siguiente oración para un ensayo formal de nivel C1:
  \emph{``The data isn't reliable.''}
\item
  ¿Qué componente obligatorio le falta a la frase: \emph{``Is very hot
  today''}?
\item
  Aplica el mecanismo de inversión a: \emph{``She is coming.''}
\end{enumerate}

👁️ Verificar análisis y respuestas

\begin{enumerate}
\def\labelenumi{\arabic{enumi}.}
\tightlist
\item
  Se debe eliminar la contracción para cumplir con el estándar
  académico: \textbf{``The data is not reliable.''}
\item
  Le falta el sujeto neutro impersonal. Lo correcto es: \textbf{``It is
  very hot today.''}
\item
  \textbf{``Is she coming?''} (El verbo pasa al frente como escudo
  interrogativo).
\end{enumerate}

\end{tcolorbox}

\chapter{CAPÍTULO 2 Los Pronombres}\label{capuxedtulo-2-los-pronombres}

\emph{Las personas del drama inglés}

\begin{center}\rule{0.5\linewidth}{0.5pt}\end{center}

\begin{quote}
\subsection{🎯 Idea Central}\label{idea-central}

En inglés \textbf{SIEMPRE} necesitas el pronombre. No puedes esconderte
como en español. \emph{``Estoy cansado''} → ❌ en inglés → ✅ \textbf{I
am tired}
\end{quote}

\begin{center}\rule{0.5\linewidth}{0.5pt}\end{center}

\section{📐 La Tabla Maestra de
Pronombres}\label{la-tabla-maestra-de-pronombres}

\begin{longtable}[]{@{}
  >{\raggedright\arraybackslash}p{(\linewidth - 10\tabcolsep) * \real{0.1343}}
  >{\raggedright\arraybackslash}p{(\linewidth - 10\tabcolsep) * \real{0.1194}}
  >{\raggedright\arraybackslash}p{(\linewidth - 10\tabcolsep) * \real{0.1194}}
  >{\raggedright\arraybackslash}p{(\linewidth - 10\tabcolsep) * \real{0.2239}}
  >{\raggedright\arraybackslash}p{(\linewidth - 10\tabcolsep) * \real{0.2388}}
  >{\raggedright\arraybackslash}p{(\linewidth - 10\tabcolsep) * \real{0.1642}}@{}}
\toprule\noalign{}
\begin{minipage}[b]{\linewidth}\raggedright
Español
\end{minipage} & \begin{minipage}[b]{\linewidth}\raggedright
Sujeto
\end{minipage} & \begin{minipage}[b]{\linewidth}\raggedright
Objeto
\end{minipage} & \begin{minipage}[b]{\linewidth}\raggedright
Posesivo Adj.
\end{minipage} & \begin{minipage}[b]{\linewidth}\raggedright
Posesivo Pron.
\end{minipage} & \begin{minipage}[b]{\linewidth}\raggedright
Reflexivo
\end{minipage} \\
\midrule\noalign{}
\endhead
\bottomrule\noalign{}
\endlastfoot
yo & \textbf{I} & me & my & mine & myself \\
tú & \textbf{you} & you & your & yours & yourself \\
él & \textbf{he} & him & his & his & himself \\
ella & \textbf{she} & her & her & hers & herself \\
ello & \textbf{it} & it & its & its & itself \\
nosotros & \textbf{we} & us & our & ours & ourselves \\
vosotros/uds & \textbf{you} & you & your & yours & yourselves \\
ellos & \textbf{they} & them & their & theirs & themselves \\
\end{longtable}

\begin{center}\rule{0.5\linewidth}{0.5pt}\end{center}

\section{💡 Truco --- Adjetivo vs Pronombre
Posesivo}\label{truco-adjetivo-vs-pronombre-posesivo}

\begin{verbatim}
ADJETIVO POSESIVO  →  acompaña a un nombre
"my car"  /  "her house"  /  "their dog"
   ↑               ↑              ↑
 siempre hay un nombre después

PRONOMBRE POSESIVO  →  va solo, sustituye al nombre
"The car is mine."  /  "The house is hers."
                ↑                      ↑
             va solo, sin nombre después
\end{verbatim}

\begin{quote}
🧠 \textbf{Regla rápida:} ¿Puedes poner un \textbf{NOMBRE} después? →
Adjetivo posesivo ¿Va \textbf{SOLO}? → Pronombre posesivo
\end{quote}

\begin{center}\rule{0.5\linewidth}{0.5pt}\end{center}

\section{✅ Ejemplos Comparados}\label{ejemplos-comparados}

\begin{verbatim}
my book     ←→   The book is mine.
your keys   ←→   The keys are yours.
his car     ←→   The car is his.
her idea    ←→   The idea is hers.
our house   ←→   The house is ours.
their dog   ←→   The dog is theirs.
\end{verbatim}

\begin{center}\rule{0.5\linewidth}{0.5pt}\end{center}

\section{⚠️ Errores Típicos}\label{errores-tuxedpicos}

\begin{verbatim}
❌  Is my book.            →  ✅  It is my book.
❌  The book is my.        →  ✅  The book is mine.
❌  I hurt me.             →  ✅  I hurt myself.
❌  She looked herself in the mirror.  →  ✅  (esto sí es correcto)
\end{verbatim}

\begin{center}\rule{0.5\linewidth}{0.5pt}\end{center}

\section{🔥 Tip para el Examen}\label{tip-para-el-examen}

\begin{quote}
\textbf{IT} hace mucho trabajo en inglés. Lo usamos para: tiempo,
distancia, temperatura y situaciones:

\begin{verbatim}
It is raining.          →  Está lloviendo.
It is far.              →  Está lejos.
It is cold.             →  Hace frío.
It is Monday.           →  Es lunes.
\end{verbatim}
\end{quote}

\chapter{CAPÍTULO 3 Present Simple}\label{capuxedtulo-3-present-simple}

\emph{El tiempo que describe tu vida}

\begin{center}\rule{0.5\linewidth}{0.5pt}\end{center}

\begin{quote}
\subsection{🎯 Idea Central}\label{idea-central-1}

Habla de \textbf{rutinas, verdades universales y hábitos.} No está
pasando AHORA --- pasa \textbf{SIEMPRE} o \textbf{REGULARMENTE.}
\end{quote}

\begin{center}\rule{0.5\linewidth}{0.5pt}\end{center}

\section{📐 La Fórmula Completa}\label{la-fuxf3rmula-completa}

\begin{verbatim}
POSITIVO:
I / You / We / They  +  VERBO BASE          →  I work
He / She / It        +  VERBO + S/ES/IES    →  He works

NEGATIVO:
I / You / We / They  +  DON'T  +  VERBO    →  I don't work
He / She / It        +  DOESN'T + VERBO    →  She doesn't work

PREGUNTA:
DO   +  I/You/We/They  +  VERBO?           →  Do you work?
DOES +  He/She/It      +  VERBO?           →  Does she work?
\end{verbatim}

\begin{center}\rule{0.5\linewidth}{0.5pt}\end{center}

\section{💡 Truco --- La ``S Viajera''}\label{truco-la-s-viajera}

\begin{verbatim}
She  workS.              ✅  (la S está en el verbo)
Does she  work?          ✅  (la S viajó a DOES)
Does she  workS?         ❌  (no puedes tener dos S)
\end{verbatim}

\begin{quote}
🧠 \textbf{Cuando usas DOES, la S ``salta'' al auxiliar.} El verbo
principal \textbf{vuelve a su forma base.}
\end{quote}

\begin{center}\rule{0.5\linewidth}{0.5pt}\end{center}

\section{💡 Truco --- ¿Cuándo añadir -S, -ES o
-IES?}\label{truco-cuuxe1ndo-auxf1adir--s--es-o--ies}

\begin{verbatim}
Verbos normales              →  +S    →  work→works, play→plays
Terminan en s,sh,ch,x,o     →  +ES   →  watch→watches, wash→washes
Consonante + Y al final      →  -Y+IES → study→studies, carry→carries
Vocal + Y al final           →  +S    →  play→plays, enjoy→enjoys
\end{verbatim}

\begin{center}\rule{0.5\linewidth}{0.5pt}\end{center}

\section{✅ Señales de Tiempo --- Time
Markers}\label{seuxf1ales-de-tiempo-time-markers}

\begin{verbatim}
always      →  siempre
usually     →  usualmente
often       →  frecuentemente
sometimes   →  a veces
rarely      →  raramente
never       →  nunca
every day   →  cada día
on Mondays  →  los lunes
in the morning → por la mañana
\end{verbatim}

\begin{center}\rule{0.5\linewidth}{0.5pt}\end{center}

\section{⚠️ Errores Típicos}\label{errores-tuxedpicos-1}

\begin{verbatim}
❌  She don't like coffee.      →  ✅  She doesn't like coffee.
❌  Does he works here?         →  ✅  Does he work here?
❌  He don't never eat meat.    →  ✅  He never eats meat.
\end{verbatim}

\begin{center}\rule{0.5\linewidth}{0.5pt}\end{center}

\section{🔥 Tip para el Examen B2}\label{tip-para-el-examen-b2}

\begin{quote}
Úsalo para \textbf{generalizaciones} en essays académicos:

\begin{verbatim}
"Research shows that people who exercise regularly live longer."
"Studies suggest that technology affects concentration."
\end{verbatim}

Suena profesional y es fácil de construir.
\end{quote}

\chapter{CAPÍTULO 4 - Present
Continuous}\label{capuxedtulo-4---present-continuous}

\emph{La foto vs.~el documental}

\begin{center}\rule{0.5\linewidth}{0.5pt}\end{center}

\begin{quote}
\subsection{🎯 Idea Central}\label{idea-central-2}

\textbf{Present Simple} = el documental de tu vida \emph{(lo que siempre
haces)} \textbf{Present Continuous} = la foto de este momento \emph{(lo
que haces AHORA)}
\end{quote}

\begin{center}\rule{0.5\linewidth}{0.5pt}\end{center}

\section{📐 La Fórmula}\label{la-fuxf3rmula}

\begin{verbatim}
POSITIVO:  Sujeto  +  AM/IS/ARE  +  VERBO-ING
NEGATIVO:  Sujeto  +  AM/IS/ARE  +  NOT  +  VERBO-ING
PREGUNTA:  AM/IS/ARE  +  Sujeto  +  VERBO-ING?

Ejemplos:
(+)  I am working.
(-)  She is not (isn't) working.
(?)  Are they working?
\end{verbatim}

\begin{center}\rule{0.5\linewidth}{0.5pt}\end{center}

\section{💡 Truco --- Las Reglas del
-ING}\label{truco-las-reglas-del--ing}

\begin{verbatim}
Verbo normal                          →  +ING       →  work  →  working
Termina en E muda                     →  quita E+ING →  make  →  making
Termina en consonante+vocal+consonante→  dobla+ING   →  run   →  running
Termina en IE                         →  IE→Y+ING    →  lie   →  lying
\end{verbatim}

\begin{center}\rule{0.5\linewidth}{0.5pt}\end{center}

\section{✅ Señales de Tiempo}\label{seuxf1ales-de-tiempo}

\begin{verbatim}
now           →  ahora
right now     →  ahora mismo
at the moment →  en este momento
Look!         →  ¡Mira! (implica algo que pasa ahora)
Listen!       →  ¡Escucha!
today         →  hoy (cuando implica este período)
this week     →  esta semana
\end{verbatim}

\begin{center}\rule{0.5\linewidth}{0.5pt}\end{center}

\section{⚠️ Los Stative Verbs --- Verbos que ODIAN el
-ING}\label{los-stative-verbs-verbos-que-odian-el--ing}

\begin{quote}
Estos verbos describen \textbf{estados mentales o emocionales}, no
acciones físicas. \textbf{NUNCA} van en continuous.
\end{quote}

\begin{verbatim}
SENTIMIENTOS:   love, hate, like, prefer, want, need, wish, adore
MENTE:          know, believe, think*, understand, remember, forget, mean
SENTIDOS:       see*, hear, smell*, taste*, feel*
POSESIÓN:       have*, own, belong, contain, include
OTROS:          seem, appear, look*, cost, weigh
\end{verbatim}

\begin{quote}
⚠️ Los marcados con * pueden usarse en continuous \textbf{con
significado diferente:}

\begin{verbatim}
I think you're right.     →  Creo que tienes razón. (opinión)
I'm thinking about you.   →  Estoy pensando en ti. (proceso activo)

I have a car.             →  Tengo un coche. (posesión)
I'm having lunch.         →  Estoy almorzando. (acción)
\end{verbatim}
\end{quote}

\begin{center}\rule{0.5\linewidth}{0.5pt}\end{center}

\section{💡 Tabla Comparativa --- Simple vs
Continuous}\label{tabla-comparativa-simple-vs-continuous}

\begin{longtable}[]{@{}ll@{}}
\toprule\noalign{}
Present Simple & Present Continuous \\
\midrule\noalign{}
\endhead
\bottomrule\noalign{}
\endlastfoot
I work in a bank. & I'm working from home today. \\
She drinks coffee every morning. & She's drinking tea right now. \\
He lives in Madrid. & He's living with his parents temporarily. \\
It rains a lot in London. & It's raining outside! \\
\end{longtable}

\begin{center}\rule{0.5\linewidth}{0.5pt}\end{center}

\section{🔥 Tip para el Examen}\label{tip-para-el-examen-1}

\begin{quote}
El Present Continuous también expresa \textbf{planes futuros} ya
confirmados:

\begin{verbatim}
"I'm meeting the director tomorrow at 10."
"They're launching the product next month."
\end{verbatim}

Esto es \textbf{clave para el writing de B2.}
\end{quote}

\part{📗 Parte 2: Nivel A2}

\chapter{CAPÍTULO 5 Past Simple}\label{capuxedtulo-5-past-simple}

\emph{El tiempo del storytelling}

\begin{center}\rule{0.5\linewidth}{0.5pt}\end{center}

\begin{quote}
\subsection{🎯 Idea Central}\label{idea-central-3}

Acción \textbf{completada} en el pasado con \textbf{tiempo definido.} La
película ya terminó. No hay conexión con el presente.
\end{quote}

\begin{center}\rule{0.5\linewidth}{0.5pt}\end{center}

\section{📐 La Fórmula}\label{la-fuxf3rmula-1}

\begin{verbatim}
POSITIVO:  Sujeto  +  VERBO PASADO (regular: +ED / irregular: forma 2)
NEGATIVO:  Sujeto  +  DIDN'T  +  VERBO BASE
PREGUNTA:  DID  +  Sujeto  +  VERBO BASE?
\end{verbatim}

\begin{center}\rule{0.5\linewidth}{0.5pt}\end{center}

\section{💡 Truco --- DID Roba el
Pasado}\label{truco-did-roba-el-pasado}

\begin{verbatim}
She worked.              →  pasado en el verbo
Did she work?            →  DID tiene el pasado, verbo vuelve a base
Did she worked?          ❌  doble pasado = error grave
She didn't work.         ✅  DIDN'T tiene el pasado
She didn't worked.       ❌
\end{verbatim}

\begin{center}\rule{0.5\linewidth}{0.5pt}\end{center}

\section{💡 Los 3 Sonidos del -ED}\label{los-3-sonidos-del--ed}

\begin{verbatim}
/t/   →  después de sonidos sordos (p,k,f,s,sh,ch):
         worked /wɜːkt/,  stopped /stɒpt/,  watched /wɒtʃt/

/d/   →  después de sonidos sonoros (b,g,v,z,m,n,l,r):
         lived /lɪvd/,  called /kɔːld/,  loved /lʌvd/

/ɪd/  →  después de T y D (necesitas una vocal para separar):
         wanted /ˈwɒntɪd/,  needed /ˈniːdɪd/,  decided /dɪˈsaɪdɪd/
\end{verbatim}

\begin{center}\rule{0.5\linewidth}{0.5pt}\end{center}

\section{📐 Los Verbos Irregulares --- Por
Familias}\label{los-verbos-irregulares-por-familias}

\begin{quote}
💡 \textbf{No los memorices uno por uno. Agrúpalos por patrones.}
\end{quote}

\textbf{Familia 1 --- Los que no cambian:}

\begin{verbatim}
cut → cut    put → put    hit → hit    let → let    set → set
🧠  "Se quedaron como estaban, como tú cuando te dicen que hay tarea"
\end{verbatim}

\textbf{Familia 2 --- Cambian la vocal del medio:}

\begin{verbatim}
run → ran       come → came      become → became
sit → sat       win → won        swim → swam
drink → drank   sing → sang      ring → rang       begin → began
drive → drove   ride → rode      write → wrote     rise → rose
\end{verbatim}

\textbf{Familia 3 --- La familia -ought / -aught:}

\begin{verbatim}
buy   → bought      think → thought     bring → brought
catch → caught      teach → taught      fight → fought
🧠  "Los que compramos, pensamos y traemos = familia BOUGHT"
\end{verbatim}

\textbf{Familia 4 --- Terminan en -eld, -ept, -elt:}

\begin{verbatim}
hold → held    keep → kept    feel → felt
sell → sold    tell → told    build → built
\end{verbatim}

\textbf{Los 20 Imprescindibles --- Memoriza estos sí o sí:}

\begin{verbatim}
be   → was/were    have → had       do   → did
go   → went        say  → said      get  → got
make → made        know → knew      take → took
see  → saw         come → came      give → gave
find → found       tell → told      feel → felt
leave → left       keep → kept      meet → met
lose → lost        choose → chose
\end{verbatim}

\begin{center}\rule{0.5\linewidth}{0.5pt}\end{center}

\section{✅ Señales de Tiempo}\label{seuxf1ales-de-tiempo-1}

\begin{verbatim}
yesterday           →  ayer
last week/month     →  la semana/el mes pasado
...ago              →  hace...  ("two years ago" = hace dos años)
in 1990             →  en 1990
when I was a child  →  cuando era niño
that day            →  ese día
\end{verbatim}

\begin{center}\rule{0.5\linewidth}{0.5pt}\end{center}

\section{⚠️ Errores Típicos}\label{errores-tuxedpicos-2}

\begin{verbatim}
❌  I did go to the cinema yesterday.     →  ✅  I went to the cinema yesterday.
❌  She didn't went home.                →  ✅  She didn't go home.
❌  Did he came early?                   →  ✅  Did he come early?
\end{verbatim}

\chapter{CAPÍTULO 6. Past
Continuous}\label{capuxedtulo-6.-past-continuous}

\emph{El escenario de la historia}

\begin{center}\rule{0.5\linewidth}{0.5pt}\end{center}

\begin{quote}
\subsection{🎯 Idea Central}\label{idea-central-4}

\textbf{Past Simple} = el rayo que cayó \emph{(acción puntual)}
\textbf{Past Continuous} = la tormenta de fondo \emph{(situación en
progreso)} Juntos cuentan la historia completa.
\end{quote}

\begin{center}\rule{0.5\linewidth}{0.5pt}\end{center}

\section{📐 La Fórmula}\label{la-fuxf3rmula-2}

\begin{verbatim}
Sujeto  +  WAS/WERE  +  VERBO-ING

I / He / She / It   →  WAS
You / We / They     →  WERE

(+)  I was sleeping.
(-)  She wasn't listening.
(?)  Were they working?
\end{verbatim}

\begin{center}\rule{0.5\linewidth}{0.5pt}\end{center}

\section{💡 La Combinación Más Usada --- WHEN y
WHILE}\label{la-combinaciuxf3n-muxe1s-usada-when-y-while}

\begin{verbatim}
━━━━━━━━━━━━━━━━━━━━━━━━━━━━━━━━━━━━━━━━━━━━
    [was sleeping — acción de fondo]
                    💥 rang
━━━━━━━━━━━━━━━━━━━━━━━━━━━━━━━━━━━━━━━━━━━━

WHEN  +  Past Simple   (la interrupción):
"I was sleeping WHEN the phone rang."

WHILE  +  Past Continuous  (la actividad de fondo):
"WHILE I was cooking, my sister arrived."
\end{verbatim}

\begin{center}\rule{0.5\linewidth}{0.5pt}\end{center}

\section{✅ Usos del Past Continuous}\label{usos-del-past-continuous}

\begin{verbatim}
1. Acción en progreso en un momento del pasado:
   "At 8pm, I was having dinner."

2. Dos acciones simultáneas en el pasado:
   "She was reading while he was cooking."

3. Descripción de fondo en una historia:
   "The sun was shining. Birds were singing. Suddenly..."
\end{verbatim}

\begin{center}\rule{0.5\linewidth}{0.5pt}\end{center}

\section{🔥 Tip para el Examen}\label{tip-para-el-examen-2}

\begin{quote}
En \textbf{narrativa y storytelling} (writing B1-B2), alterna Past
Simple y Past Continuous:

\begin{verbatim}
"It was a dark night. The wind was blowing and the rain
was hitting the windows when suddenly the lights went out."
\end{verbatim}

Esto demuestra dominio del tiempo narrativo.
\end{quote}

\chapter{CAPÍTULO 7 Future Forms}\label{capuxedtulo-7-future-forms}

\emph{Tres formas de hablar del futuro --- y cuándo usar cada una}

\begin{center}\rule{0.5\linewidth}{0.5pt}\end{center}

\begin{quote}
\subsection{🎯 Idea Central}\label{idea-central-5}

El inglés tiene \textbf{múltiples futuros} y cada uno dice algo
diferente sobre tu \textbf{relación} con ese evento futuro. No son
intercambiables. Cada uno tiene su momento.
\end{quote}

\begin{center}\rule{0.5\linewidth}{0.5pt}\end{center}

\section{📐 Las Tres Formas
Principales}\label{las-tres-formas-principales}

\begin{center}\rule{0.5\linewidth}{0.5pt}\end{center}

\subsection{🔮 WILL --- El futuro
espontáneo}\label{will-el-futuro-espontuxe1neo}

\begin{verbatim}
Sujeto  +  WILL  +  VERBO BASE
Negativo:  won't (will not)
\end{verbatim}

\textbf{Cuándo usarlo:}

\begin{verbatim}
✅  Decisiones espontáneas (decides en el momento):
    "I'll have the pizza."  (decides en el restaurante, ahora mismo)

✅  Predicciones sin evidencia concreta (opinión):
    "I think it will rain tomorrow."

✅  Promesas y ofrecimientos:
    "I'll help you."  /  "I'll call you later."

✅  Verdades y hechos futuros:
    "She'll be 30 next year."
\end{verbatim}

\begin{center}\rule{0.5\linewidth}{0.5pt}\end{center}

\subsection{📅 GOING TO --- El futuro
planificado}\label{going-to-el-futuro-planificado}

\begin{verbatim}
Sujeto  +  AM/IS/ARE  +  GOING TO  +  VERBO BASE
\end{verbatim}

\textbf{Cuándo usarlo:}

\begin{verbatim}
✅  Planes ya decididos antes de hablar:
    "I'm going to study medicine."  (ya lo tenías decidido)

✅  Predicciones con EVIDENCIA VISIBLE:
    "Look at those clouds! It's going to rain."
    🧠 Si PUEDES VER la evidencia → going to
\end{verbatim}

\begin{center}\rule{0.5\linewidth}{0.5pt}\end{center}

\subsection{🗓️ PRESENT CONTINUOUS --- El futuro en
agenda}\label{present-continuous-el-futuro-en-agenda}

\begin{verbatim}
Sujeto  +  AM/IS/ARE  +  VERBO-ING
\end{verbatim}

\textbf{Cuándo usarlo:}

\begin{verbatim}
✅  Planes con hora y lugar ya fijados:
    "I'm meeting John at 6."       (está en el calendario)
    "They're getting married in June."  (todo organizado)
    "We're flying to Rome on Friday."
\end{verbatim}

\begin{center}\rule{0.5\linewidth}{0.5pt}\end{center}

\section{💡 Tabla Comparativa --- El Truco
Visual}\label{tabla-comparativa-el-truco-visual}

\begin{longtable}[]{@{}
  >{\raggedright\arraybackslash}p{(\linewidth - 2\tabcolsep) * \real{0.4231}}
  >{\raggedright\arraybackslash}p{(\linewidth - 2\tabcolsep) * \real{0.5769}}@{}}
\toprule\noalign{}
\begin{minipage}[b]{\linewidth}\raggedright
Situación
\end{minipage} & \begin{minipage}[b]{\linewidth}\raggedright
Forma correcta
\end{minipage} \\
\midrule\noalign{}
\endhead
\bottomrule\noalign{}
\endlastfoot
El teléfono suena. Decides: \emph{``¡Yo lo cojo!''} & \textbf{WILL} \\
Ya decidiste: \emph{``Este verano viajo a Londres''} & \textbf{GOING
TO} \\
Tienes cita: \emph{``Ceno con María el viernes a las 8''} &
\textbf{PRESENT CONTINUOUS} \\
Ves nubes negras en el cielo & \textbf{GOING TO} \\
Te piden que predices el futuro sin datos & \textbf{WILL} \\
\end{longtable}

\begin{center}\rule{0.5\linewidth}{0.5pt}\end{center}

\section{⚠️ Errores Típicos}\label{errores-tuxedpicos-3}

\begin{verbatim}
❌  If it will rain, I'll stay home.
✅  If it rains, I'll stay home.
    (después de IF nunca uses WILL — ver Capítulo 14)

❌  I going to travel next year.
✅  I'm going to travel next year.

❌  I will meet her at 8 tonight. (tienes cita confirmada)
✅  I'm meeting her at 8 tonight.
\end{verbatim}

\part{📙 Parte 3: Nivel B1}

\chapter{CAPÍTULO 8 Present
Perfect}\label{capuxedtulo-8-present-perfect}

\emph{El tiempo que más confunde a los hispanohablantes}

\begin{center}\rule{0.5\linewidth}{0.5pt}\end{center}

\begin{quote}
\subsection{🎯 Idea Central}\label{idea-central-6}

Conecta el \textbf{pasado con el presente.} No importa exactamente
cuándo pasó. Lo que importa es que \textbf{tiene relevancia AHORA.}
\end{quote}

\begin{center}\rule{0.5\linewidth}{0.5pt}\end{center}

\begin{quote}
\subsection{🧠 El Gran Problema
Español}\label{el-gran-problema-espauxf1ol}

En español decimos: - \emph{``He comido''} → Present Perfect -
\emph{``Ayer comí''} → también podría traducirse igual

El inglés \textbf{separa estas dos cosas radicalmente:} - Sin tiempo
definido → \textbf{Present Perfect} - Con tiempo definido del pasado →
\textbf{Past Simple}

\textbf{Ahí está todo el lío. Y ahora ya lo sabes.}
\end{quote}

\begin{center}\rule{0.5\linewidth}{0.5pt}\end{center}

\section{📐 La Fórmula}\label{la-fuxf3rmula-3}

\begin{verbatim}
Sujeto  +  HAVE / HAS  +  PARTICIPIO PASADO

I / You / We / They   →  HAVE
He / She / It         →  HAS

(+)  I have eaten.          /  She has left.
(-)  I haven't eaten.       /  She hasn't left.
(?)  Have you eaten?        /  Has she left?
\end{verbatim}

\begin{center}\rule{0.5\linewidth}{0.5pt}\end{center}

\section{💡 Los Participios Irregulares Más
Importantes}\label{los-participios-irregulares-muxe1s-importantes}

\begin{verbatim}
be    → been        go    → gone        do    → done
have  → had         see   → seen        take  → taken
give  → given       write → written     eat   → eaten
fall  → fallen      speak → spoken      break → broken
buy   → bought      think → thought     bring → brought
make  → made        come  → come        run   → run
know  → known       grow  → grown       show  → shown
\end{verbatim}

\begin{center}\rule{0.5\linewidth}{0.5pt}\end{center}

\section{✅ Las Señales de Tiempo del Present
Perfect}\label{las-seuxf1ales-de-tiempo-del-present-perfect}

\begin{verbatim}
just       →  acabar de:        "I've just eaten."
ever       →  alguna vez:       "Have you ever been to Japan?"
never      →  nunca:            "I've never tried sushi."
already    →  ya:               "She's already left."
yet        →  todavía / ya:     "I haven't finished yet." / "Have you eaten yet?"
since      →  desde:            "I've lived here since 2010."
for        →  durante:          "I've lived here for 5 years."
recently   →  recientemente:    "I've recently started running."
lately     →  últimamente:      "Have you been okay lately?"
so far     →  hasta ahora:      "So far, I've read 3 chapters."
\end{verbatim}

\begin{center}\rule{0.5\linewidth}{0.5pt}\end{center}

\section{💡 Truco --- SINCE vs FOR}\label{truco-since-vs-for}

\begin{verbatim}
SINCE  →  punto específico en el tiempo (cuándo empezó)
FOR    →  duración (cuánto tiempo lleva)

SINCE:   since 2010 / since Monday / since I was a child / since we met
FOR:     for 5 years / for two months / for a long time / for ages

🧠  SINCE = Since a Specific moment
🧠  FOR   = FOR a duration (cantidad de tiempo)
\end{verbatim}

\begin{center}\rule{0.5\linewidth}{0.5pt}\end{center}

\section{💡 La Regla de Oro --- Perfect vs
Simple}\label{la-regla-de-oro-perfect-vs-simple}

\begin{verbatim}
┌─────────────────────────────────────────────────────────┐
│  ¿Hay tiempo DEFINIDO del pasado en la frase?           │
│  (yesterday, last week, in 2010, ago, when I was...)    │
│                                                         │
│  SÍ  →  PAST SIMPLE                                     │
│  NO  →  PRESENT PERFECT                                 │
└─────────────────────────────────────────────────────────┘
\end{verbatim}

\begin{center}\rule{0.5\linewidth}{0.5pt}\end{center}

\section{✅ Tabla Comparativa --- Perfect vs
Simple}\label{tabla-comparativa-perfect-vs-simple}

\begin{longtable}[]{@{}
  >{\raggedright\arraybackslash}p{(\linewidth - 2\tabcolsep) * \real{0.5517}}
  >{\raggedright\arraybackslash}p{(\linewidth - 2\tabcolsep) * \real{0.4483}}@{}}
\toprule\noalign{}
\begin{minipage}[b]{\linewidth}\raggedright
Present Perfect
\end{minipage} & \begin{minipage}[b]{\linewidth}\raggedright
Past Simple
\end{minipage} \\
\midrule\noalign{}
\endhead
\bottomrule\noalign{}
\endlastfoot
I've lost my keys. \emph{(siguen perdidas)} & I lost my keys
yesterday. \\
She's broken her arm. & She broke her arm in 2020. \\
Have you ever eaten Thai food? & Did you eat Thai food last night? \\
He's written three novels. \emph{(logro)} & He wrote his first novel in
2005. \\
\end{longtable}

\begin{center}\rule{0.5\linewidth}{0.5pt}\end{center}

\section{⚠️ Error Crítico --- El Más
Común}\label{error-cruxedtico-el-muxe1s-comuxfan}

\begin{verbatim}
❌  I have seen her yesterday.
✅  I saw her yesterday.

❌  Did you ever try sushi?
✅  Have you ever tried sushi?

❌  I have gone to Paris in 2015.
✅  I went to Paris in 2015.
\end{verbatim}

\chapter{CAPÍTULO 9 Present Perfect
Continuous}\label{capuxedtulo-9-present-perfect-continuous}

\emph{``Llevo haciendo esto un tiempo\ldots{} y se nota''}

\begin{center}\rule{0.5\linewidth}{0.5pt}\end{center}

\begin{quote}
\subsection{🎯 Idea Central}\label{idea-central-7}

Acción que \textbf{empezó en el pasado, continúa ahora} y tiene un
\textbf{resultado o efecto visible en el presente.}
\end{quote}

\begin{center}\rule{0.5\linewidth}{0.5pt}\end{center}

\section{📐 La Fórmula}\label{la-fuxf3rmula-4}

\begin{verbatim}
Sujeto  +  HAVE/HAS  +  BEEN  +  VERBO-ING

(+)  I have been studying for 3 hours.
(-)  She hasn't been sleeping well.
(?)  Have they been waiting long?
\end{verbatim}

\begin{center}\rule{0.5\linewidth}{0.5pt}\end{center}

\section{✅ Ejemplos con Contexto}\label{ejemplos-con-contexto}

\begin{verbatim}
"I've been studying for 3 hours."
→ Llevo 3 horas estudiando (y todavía sigo / o acabo de parar agotado)

"She's been crying."
→ Ha estado llorando (se le nota: ojos rojos, cara hinchada)

"They've been working all day."
→ Han estado trabajando todo el día (están exhaustos ahora)

"It's been raining."
→ Ha estado lloviendo (el suelo está mojado)
\end{verbatim}

\begin{center}\rule{0.5\linewidth}{0.5pt}\end{center}

\section{💡 Truco --- Perfect vs Perfect
Continuous}\label{truco-perfect-vs-perfect-continuous}

\begin{verbatim}
PRESENT PERFECT           →  resultado / logro / cantidad
"I've written 3 emails."  →  (están escritos, importa el resultado)
"She's read the report."  →  (lo terminó, está hecho)

PRESENT PERFECT CONTINUOUS  →  proceso / duración / actividad en sí
"I've been writing emails all morning."
→ (importa que llevo tiempo haciéndolo, quizás no terminé)
"She's been reading the report."
→ (está en proceso, puede que no haya terminado)
\end{verbatim}

\chapter{CAPÍTULO 10 Past Perfect}\label{capuxedtulo-10-past-perfect}

\emph{El pasado del pasado}

\begin{center}\rule{0.5\linewidth}{0.5pt}\end{center}

\begin{quote}
\subsection{🎯 Idea Central}\label{idea-central-8}

Cuando ya estás hablando del pasado y necesitas referirte a algo que
pasó \textbf{AÚN ANTES.} Es el \textbf{flashback} gramatical del inglés.
\end{quote}

\begin{center}\rule{0.5\linewidth}{0.5pt}\end{center}

\section{📐 La Fórmula}\label{la-fuxf3rmula-5}

\begin{verbatim}
Sujeto  +  HAD  +  PARTICIPIO PASADO
(HAD es igual para todos los sujetos — nunca cambia)

(+)  She had left before I arrived.
(-)  He hadn't eaten when we called.
(?)  Had they finished when you got there?
\end{verbatim}

\begin{center}\rule{0.5\linewidth}{0.5pt}\end{center}

\section{💡 Línea de Tiempo Visual}\label{luxednea-de-tiempo-visual}

\begin{verbatim}
PASADO  ←──────────────────────────────────────→  PRESENTE

    HAD FINISHED          ARRIVED
           │                 │
    ───────●─────────────────●──────────────────▶
    (pasado del          (pasado)
      pasado)

"When I arrived, the party had already finished."
 → Primero terminó la fiesta, DESPUÉS llegué yo.
\end{verbatim}

\begin{center}\rule{0.5\linewidth}{0.5pt}\end{center}

\section{✅ Ejemplos en Contexto}\label{ejemplos-en-contexto}

\begin{verbatim}
"She was tired because she hadn't slept."
→ Primero no durmió → después estaba cansada.

"He recognized her because he had seen her before."
→ Primero la vio → después la reconoció.

"By the time we arrived, they had already left."
→ Se fueron antes de que llegáramos.

"I realized I had forgotten my wallet."
→ Me di cuenta de que lo había olvidado (antes).
\end{verbatim}

\begin{center}\rule{0.5\linewidth}{0.5pt}\end{center}

\section{🔥 Tip para el Examen}\label{tip-para-el-examen-3}

\begin{quote}
El Past Perfect es \textbf{CLAVE} en: - \textbf{Reported Speech}
\emph{(ver Capítulo 13)} - \textbf{Third Conditional} \emph{(ver
Capítulo 14)} - \textbf{Narrativa avanzada} en writing

Si dominas el Past Perfect, dominas el B2.
\end{quote}

\chapter{CAPÍTULO 11 Modal Verbs}\label{capuxedtulo-11-modal-verbs}

\emph{Los verbos con personalidad propia}

\begin{center}\rule{0.5\linewidth}{0.5pt}\end{center}

\begin{quote}
\subsection{🎯 Idea Central}\label{idea-central-9}

Los modales no expresan acciones. Expresan tu \textbf{ACTITUD} hacia la
acción: obligación, posibilidad, permiso, capacidad, consejo.
\end{quote}

\begin{center}\rule{0.5\linewidth}{0.5pt}\end{center}

\section{📐 Reglas Universales de los
Modales}\label{reglas-universales-de-los-modales}

\begin{verbatim}
1.  NUNCA llevan -S en tercera persona:
    ❌  He cans swim.    →  ✅  He can swim.

2.  NUNCA van con TO (excepto: ought to, have to, used to):
    ❌  She must to go.  →  ✅  She must go.

3.  SIEMPRE van seguidos de VERBO BASE:
    ❌  I should studied. →  ✅  I should study.

4.  Forman negación directamente con NOT:
    ❌  I don't can swim. →  ✅  I can't swim.
\end{verbatim}

\begin{center}\rule{0.5\linewidth}{0.5pt}\end{center}

\section{📐 La Familia Completa de
Modales}\label{la-familia-completa-de-modales}

\begin{center}\rule{0.5\linewidth}{0.5pt}\end{center}

\subsection{🔵 CAN / COULD --- Habilidad y
Permiso}\label{can-could-habilidad-y-permiso}

\begin{verbatim}
CAN (presente / informal):
"I can swim."                →  Sé nadar.
"Can I open the window?"     →  ¿Puedo abrir la ventana? (informal)

COULD (pasado / formal / tentativo):
"I could swim when I was 5." →  Sabía nadar cuando tenía 5 años.
"Could I use your phone?"    →  ¿Podría usar tu teléfono? (formal)
"Could you help me?"         →  ¿Podrías ayudarme? (petición educada)
\end{verbatim}

\begin{center}\rule{0.5\linewidth}{0.5pt}\end{center}

\subsection{🟡 MUST / HAVE TO ---
Obligación}\label{must-have-to-obligaciuxf3n}

\begin{verbatim}
MUST  →  obligación INTERNA (viene de ti mismo)
"I must study."              →  Debo estudiar. (me lo digo a mí mismo)
"You must try this!"         →  ¡Tienes que probarlo! (recomendación fuerte)

HAVE TO  →  obligación EXTERNA (reglas, otras personas)
"I have to wear a uniform."  →  Tengo que llevar uniforme. (reglas del trabajo)
"She has to be home by 11."  →  Tiene que estar en casa. (sus padres lo dicen)
\end{verbatim}

\begin{quote}
🧠 \textbf{Truco definitivo:} \textbf{MUST} = viene de \textbf{MÍ}
\textbf{HAVE TO} = viene de \textbf{FUERA}
\end{quote}

\begin{center}\rule{0.5\linewidth}{0.5pt}\end{center}

\subsection{🔴 MUSTN'T vs DON'T HAVE TO --- ¡Los más
confundidos!}\label{mustnt-vs-dont-have-to-los-muxe1s-confundidos}

\begin{verbatim}
┌────────────────────────────────────────────────────────┐
│  MUSTN'T  =  PROHIBICIÓN  (NO debes hacer esto)        │
│  "You mustn't smoke here."  →  Está PROHIBIDO fumar.   │
│                                                         │
│  DON'T HAVE TO  =  NO ES NECESARIO (pero puedes)       │
│  "You don't have to come."  →  No tienes que venir.    │
│                                (es opcional, decides tú)│
└────────────────────────────────────────────────────────┘
\end{verbatim}

\begin{quote}
⚠️ \textbf{Este error puede ser catastrófico en contexto real:}

\begin{verbatim}
Situación: la entrada es gratuita.
❌  "You mustn't pay."     →  suena como que está PROHIBIDO pagar
✅  "You don't have to pay."  →  no es necesario pagar
\end{verbatim}
\end{quote}

\begin{center}\rule{0.5\linewidth}{0.5pt}\end{center}

\subsection{🟢 SHOULD / OUGHT TO --- Consejo y
Recomendación}\label{should-ought-to-consejo-y-recomendaciuxf3n}

\begin{verbatim}
SHOULD  →  consejo, recomendación, opinión
"You should see a doctor."    →  Deberías ver a un médico.
"Should I tell her?"          →  ¿Debería decírselo?

OUGHT TO  →  igual que should, ligeramente más formal
"You ought to apologize."     →  Deberías disculparte.
"They ought to arrive soon."  →  Deberían llegar pronto.
\end{verbatim}

\begin{center}\rule{0.5\linewidth}{0.5pt}\end{center}

\subsection{⚪ MAY / MIGHT --- Posibilidad}\label{may-might-posibilidad}

\begin{verbatim}
MAY    →  posibilidad más alta (~50% o más)
"It may rain today."          →  Puede que llueva hoy.
"May I come in?"              →  ¿Puedo pasar? (permiso MUY formal)

MIGHT  →  posibilidad más baja (~30% o menos)
"I might go to the party."    →  Igual voy a la fiesta (no estoy seguro).
"She might be late."          →  Puede que llegue tarde.
\end{verbatim}

\begin{quote}
🧠 \textbf{Truco:} MIGHT = \textbf{Maybe I will, maybe I won't}
\end{quote}

\begin{center}\rule{0.5\linewidth}{0.5pt}\end{center}

\section{📐 Tabla Resumen de Modales}\label{tabla-resumen-de-modales}

\begin{longtable}[]{@{}
  >{\raggedright\arraybackslash}p{(\linewidth - 4\tabcolsep) * \real{0.2800}}
  >{\raggedright\arraybackslash}p{(\linewidth - 4\tabcolsep) * \real{0.3600}}
  >{\raggedright\arraybackslash}p{(\linewidth - 4\tabcolsep) * \real{0.3600}}@{}}
\toprule\noalign{}
\begin{minipage}[b]{\linewidth}\raggedright
Modal
\end{minipage} & \begin{minipage}[b]{\linewidth}\raggedright
Función
\end{minipage} & \begin{minipage}[b]{\linewidth}\raggedright
Ejemplo
\end{minipage} \\
\midrule\noalign{}
\endhead
\bottomrule\noalign{}
\endlastfoot
\textbf{can} & habilidad / permiso informal & I can drive. \\
\textbf{could} & habilidad pasada / petición formal & Could you help? \\
\textbf{must} & obligación interna / deducción fuerte & You must be
tired. \\
\textbf{have to} & obligación externa & I have to work. \\
\textbf{mustn't} & prohibición absoluta & You mustn't lie. \\
\textbf{don't have to} & no es necesario & You don't have to pay. \\
\textbf{should} & consejo / recomendación & You should rest. \\
\textbf{ought to} & consejo formal & You ought to apologize. \\
\textbf{may} & posibilidad alta / permiso formal & It may snow. \\
\textbf{might} & posibilidad baja & I might come. \\
\textbf{would} & condicional / petición educada & Would you
like\ldots? \\
\textbf{shall} & sugerencia / ofrecimiento (formal) & Shall we go? \\
\end{longtable}

\begin{center}\rule{0.5\linewidth}{0.5pt}\end{center}

\section{🔥 Modales Perfectos --- Nivel B2
Puro}\label{modales-perfectos-nivel-b2-puro}

\begin{quote}
\textbf{Estructura:} MODAL + HAVE + PARTICIPIO PASADO Usados para hacer
\textbf{deducciones, críticas y especulaciones sobre el pasado.}
\end{quote}

\begin{verbatim}
MUST HAVE + participio     →  deducción casi segura sobre el pasado
"She must have forgotten." →  Debe haber olvidado. (estoy casi seguro)

CAN'T HAVE + participio    →  imposibilidad en el pasado
"He can't have said that." →  No puede haber dicho eso. (imposible)

SHOULD HAVE + participio   →  crítica o lamento sobre el pasado
"I should have studied."   →  Debería haber estudiado. (pero no lo hice)

COULD HAVE + participio    →  posibilidad no aprovechada
"You could have told me!"  →  ¡Podrías haberme dicho! (pero no lo hiciste)

MIGHT HAVE + participio    →  posibilidad en el pasado
"She might have missed it."→  Puede que lo haya perdido.

NEEDN'T HAVE + participio  →  hiciste algo que no era necesario
"You needn't have cooked." →  No hacía falta que cocinaras.
\end{verbatim}

\chapter{CAPÍTULO 12 Passive Voice}\label{capuxedtulo-12-passive-voice}

\emph{Cuando el QUÉ importa más que el QUIÉN}

\begin{center}\rule{0.5\linewidth}{0.5pt}\end{center}

\begin{quote}
\subsection{🎯 Idea Central}\label{idea-central-10}

Usamos la pasiva cuando: 1. No sabemos quién hizo la acción 2. No
importa quién la hizo 3. Es obvio quién la hizo 4. Queremos sonar más
\textbf{formal o académico}
\end{quote}

\begin{center}\rule{0.5\linewidth}{0.5pt}\end{center}

\section{📐 La Fórmula Maestra}\label{la-fuxf3rmula-maestra}

\begin{verbatim}
Sujeto  +  TO BE (tiempo correcto)  +  PARTICIPIO PASADO
                        ↑
            Este es el que cambia según el tiempo verbal
\end{verbatim}

\begin{center}\rule{0.5\linewidth}{0.5pt}\end{center}

\section{📐 Tabla de Transformación por
Tiempos}\label{tabla-de-transformaciuxf3n-por-tiempos}

\begin{longtable}[]{@{}
  >{\raggedright\arraybackslash}p{(\linewidth - 4\tabcolsep) * \real{0.2581}}
  >{\raggedright\arraybackslash}p{(\linewidth - 4\tabcolsep) * \real{0.3548}}
  >{\raggedright\arraybackslash}p{(\linewidth - 4\tabcolsep) * \real{0.3871}}@{}}
\toprule\noalign{}
\begin{minipage}[b]{\linewidth}\raggedright
Tiempo
\end{minipage} & \begin{minipage}[b]{\linewidth}\raggedright
Voz Activa
\end{minipage} & \begin{minipage}[b]{\linewidth}\raggedright
Voz Pasiva
\end{minipage} \\
\midrule\noalign{}
\endhead
\bottomrule\noalign{}
\endlastfoot
Present Simple & They make cars here. & Cars \textbf{are made} here. \\
Past Simple & They built this in 1900. & This \textbf{was built} in
1900. \\
Present Continuous & They are filming it. & It \textbf{is being
filmed.} \\
Past Continuous & They were repairing it. & It \textbf{was being
repaired.} \\
Present Perfect & They have arrested him. & He \textbf{has been
arrested.} \\
Past Perfect & They had warned us. & We \textbf{had been warned.} \\
Future (will) & They will announce it. & It \textbf{will be
announced.} \\
Modal & They must fix it. & It \textbf{must be fixed.} \\
Infinitive & They want to sell it. & It \textbf{is to be sold.} \\
\end{longtable}

\begin{center}\rule{0.5\linewidth}{0.5pt}\end{center}

\section{💡 Los 3 Pasos de
Transformación}\label{los-3-pasos-de-transformaciuxf3n}

\begin{verbatim}
ACTIVA:   The chef    cooks    the meal    every day.
               ↓         ↓         ↓
PASIVA:   The meal   IS COOKED  (by the chef)  every day.

Paso 1:  El OBJETO de la activa → SUJETO de la pasiva
Paso 2:  TO BE se conjuga en el mismo tiempo
Paso 3:  El verbo → PARTICIPIO PASADO
         (el agente "by + persona" es opcional)
\end{verbatim}

\begin{center}\rule{0.5\linewidth}{0.5pt}\end{center}

\section{✅ Cuándo Usar o No el ``By +
Agente''}\label{cuuxe1ndo-usar-o-no-el-by-agente}

\begin{verbatim}
Se OMITE cuando:
- Es obvio:   "He was arrested."  (por la policía, obvio)
- No se sabe: "The window was broken."  (no sabemos quién)

Se INCLUYE cuando:
- Es importante saber quién:
  "The theory was developed by Einstein."
  "The novel was written by García Márquez."
\end{verbatim}

\begin{center}\rule{0.5\linewidth}{0.5pt}\end{center}

\section{⚠️ Errores Frecuentes}\label{errores-frecuentes}

\begin{verbatim}
❌  The car was broke.          →  ✅  The car was broken.
❌  It was builded in 1800.     →  ✅  It was built in 1800.
❌  The letter has been writed. →  ✅  The letter has been written.
❌  It was been stolen.         →  ✅  It was stolen. / It has been stolen.
\end{verbatim}

\begin{center}\rule{0.5\linewidth}{0.5pt}\end{center}

\section{🔥 Tip para el Examen B2}\label{tip-para-el-examen-b2-1}

\begin{quote}
La pasiva es \textbf{muy frecuente} en writing formal y académico.
Aprende estas estructuras de memoria:

\begin{verbatim}
"It is believed that..."          →  Se cree que...
"It has been proven that..."      →  Se ha demostrado que...
"It is widely accepted that..."   →  Es ampliamente aceptado que...
"Research has been conducted..."  →  Se ha realizado una investigación...
"The results were published in..."→  Los resultados fueron publicados en...
\end{verbatim}
\end{quote}

\chapter{CAPÍTULO 13 Reported
Speech}\label{capuxedtulo-13-reported-speech}

\emph{``Me dijeron que\ldots{}'' --- El arte de contar lo que dijo
alguien}

\begin{center}\rule{0.5\linewidth}{0.5pt}\end{center}

\begin{quote}
\subsection{🎯 Idea Central}\label{idea-central-11}

Cuando reportas lo que alguien dijo, el tiempo verbal \textbf{retrocede
un paso} en el pasado. Es como si el tiempo ``se alejara'' un nivel.
\end{quote}

\begin{center}\rule{0.5\linewidth}{0.5pt}\end{center}

\section{💡 La Escalera del Tiempo ---
Backshift}\label{la-escalera-del-tiempo-backshift}

\begin{verbatim}
╔══════════════════════════════════════════════════════════╗
║  DIRECT SPEECH          →    REPORTED SPEECH            ║
╠══════════════════════════════════════════════════════════╣
║  Present Simple         →    Past Simple                 ║
║  "I work here."              She said she worked there.  ║
╠══════════════════════════════════════════════════════════╣
║  Present Continuous     →    Past Continuous             ║
║  "I am working."             She said she was working.   ║
╠══════════════════════════════════════════════════════════╣
║  Past Simple            →    Past Perfect                ║
║  "I worked hard."            She said she had worked.    ║
╠══════════════════════════════════════════════════════════╣
║  Present Perfect        →    Past Perfect                ║
║  "I have seen it."           She said she had seen it.   ║
╠══════════════════════════════════════════════════════════╣
║  Will                   →    Would                       ║
║  "I will call."              She said she would call.    ║
╠══════════════════════════════════════════════════════════╣
║  Can                    →    Could                       ║
║  "I can help."               She said she could help.    ║
╠══════════════════════════════════════════════════════════╣
║  May                    →    Might                       ║
║  "I may come."               She said she might come.    ║
╠══════════════════════════════════════════════════════════╣
║  Must                   →    Had to                      ║
║  "I must go."                She said she had to go.     ║
╚══════════════════════════════════════════════════════════╝
\end{verbatim}

\begin{center}\rule{0.5\linewidth}{0.5pt}\end{center}

\section{💡 También Cambian las Expresiones de Tiempo y
Lugar}\label{tambiuxe9n-cambian-las-expresiones-de-tiempo-y-lugar}

\begin{verbatim}
here         →  there
now          →  then
today        →  that day
tomorrow     →  the next day / the following day
yesterday    →  the day before / the previous day
this         →  that
these        →  those
...ago       →  ...before
next week    →  the following week
last week    →  the previous week
\end{verbatim}

\begin{center}\rule{0.5\linewidth}{0.5pt}\end{center}

\section{📐 Reported Questions --- Cuidado con el
Orden}\label{reported-questions-cuidado-con-el-orden}

\begin{verbatim}
Wh- Questions:
Direct:   "Where do you live?" he asked.
Reported: He asked where I lived.
                              ↑
                    orden afirmativo (sin inversión)
                    + tiempo atrás

Yes/No Questions  →  IF o WHETHER:
Direct:   "Are you coming?" she asked.
Reported: She asked if / whether I was coming.
\end{verbatim}

\begin{center}\rule{0.5\linewidth}{0.5pt}\end{center}

\section{📐 Reported Orders e
Instrucciones}\label{reported-orders-e-instrucciones}

\begin{verbatim}
Órdenes:
"Open the window!"     →  He told me TO OPEN the window.
"Don't touch that!"    →  She told me NOT TO TOUCH that.

Peticiones:
"Please help me."      →  She asked me TO HELP her.

Consejo:
"You should rest."     →  He advised me TO REST.
\end{verbatim}

\begin{center}\rule{0.5\linewidth}{0.5pt}\end{center}

\section{⚠️ Errores Típicos}\label{errores-tuxedpicos-4}

\begin{verbatim}
❌  He said me that he was tired.   →  ✅  He told me that he was tired.
❌  She said that she is tired.     →  ✅  She said that she was tired.
❌  He asked where did I live.      →  ✅  He asked where I lived.
\end{verbatim}

\begin{quote}
🧠 \textbf{SAY vs TELL:}

\begin{verbatim}
SAY  →  sin objeto directo de persona:  He said (that)...
TELL →  con objeto directo de persona:  He told ME (that)...
\end{verbatim}
\end{quote}

\part{📕 Parte 4: Nivel B2}

\chapter{CAPÍTULO 14 -
Conditionals}\label{capuxedtulo-14---conditionals}

\emph{El capítulo más temido --- Desmitificado para siempre}

\begin{center}\rule{0.5\linewidth}{0.5pt}\end{center}

\begin{quote}
\subsection{🎯 Idea Central}\label{idea-central-12}

Los condicionales hablan de \textbf{condiciones y sus consecuencias.} La
clave está en entender qué tan \textbf{REAL o POSIBLE} es la situación.
Una vez entiendes eso, las fórmulas tienen sentido lógico.
\end{quote}

\begin{center}\rule{0.5\linewidth}{0.5pt}\end{center}

\section{📐 Los Cuatro Condicionales}\label{los-cuatro-condicionales}

\begin{center}\rule{0.5\linewidth}{0.5pt}\end{center}

\subsection{🟢 ZERO CONDITIONAL --- Las verdades
universales}\label{zero-conditional-las-verdades-universales}

\begin{verbatim}
IF  +  Present Simple,  Present Simple

"If you heat water to 100°C, it boils."
"If I don't sleep, I feel terrible."
"Plants die if they don't get water."

Cuándo usarlo:
✅  Hechos científicos y verdades universales
✅  Instrucciones y procesos
✅  Hábitos y rutinas (IF = WHEN, son intercambiables aquí)
\end{verbatim}

\begin{center}\rule{0.5\linewidth}{0.5pt}\end{center}

\subsection{🔵 FIRST CONDITIONAL --- Lo que puede pasar de
verdad}\label{first-conditional-lo-que-puede-pasar-de-verdad}

\begin{verbatim}
IF  +  Present Simple,  WILL  +  Verbo Base

"If it rains, I will take an umbrella."
"If you study hard, you will pass."
"I'll call you if I'm late."

Cuándo usarlo:
✅  Situación REAL y POSIBLE en el futuro
✅  Hay una probabilidad genuina de que ocurra
⚠️  NUNCA uses WILL después de IF: "If it will rain" ❌
\end{verbatim}

\begin{center}\rule{0.5\linewidth}{0.5pt}\end{center}

\subsection{🟡 SECOND CONDITIONAL --- La fantasía del
presente/futuro}\label{second-conditional-la-fantasuxeda-del-presentefuturo}

\begin{verbatim}
IF  +  Past Simple,  WOULD  +  Verbo Base

"If I had a million dollars, I would travel the world."
"If I were you, I would apologize."
"What would you do if you won the lottery?"

Cuándo usarlo:
✅  Situación IRREAL o muy improbable AHORA o en el FUTURO
✅  Estás imaginando una realidad alternativa
✅  Dar consejos de forma indirecta
\end{verbatim}

\begin{quote}
💡 \textbf{Truco especial --- IF I WERE:}

\begin{verbatim}
En inglés formal y B2:  "If I were you..."  ✅
En inglés informal:     "If I was you..."   ✅ (aceptable)
En B2 escrito:          usa WERE siempre    🔥
\end{verbatim}
\end{quote}

\begin{center}\rule{0.5\linewidth}{0.5pt}\end{center}

\subsection{🔴 THIRD CONDITIONAL --- El arrepentimiento del
pasado}\label{third-conditional-el-arrepentimiento-del-pasado}

\begin{verbatim}
IF  +  Past Perfect,  WOULD HAVE  +  Participio

"If I had studied, I would have passed the exam."
"If she had called, I would have answered."
"We wouldn't have been late if we had left earlier."

Cuándo usarlo:
✅  Situación IRREAL en el PASADO — ya no se puede cambiar
✅  Lamentos, arrepentimientos, críticas
✅  Especular sobre situaciones pasadas alternativas
\end{verbatim}

\begin{center}\rule{0.5\linewidth}{0.5pt}\end{center}

\section{💡 El Mapa Visual de los
Condicionales}\label{el-mapa-visual-de-los-condicionales}

\begin{verbatim}
REAL ◄─────────────────────────────────► IRREAL

   0              1              2              3
SIEMPRE        POSIBLE       IMPROBABLE     IMPOSIBLE
                              /irreal       (ya pasó)

agua→vapor   llueve→       millonario→    estudié más→
             paraguas       viajo          aprobé
\end{verbatim}

\begin{center}\rule{0.5\linewidth}{0.5pt}\end{center}

\section{🔥 Mixed Conditionals --- Solo
B2}\label{mixed-conditionals-solo-b2}

\begin{quote}
Mezclas el tiempo de la \textbf{condición} y el \textbf{resultado:}
\end{quote}

\begin{verbatim}
PASADO IRREAL  →  consecuencia en el PRESENTE:
IF + Past Perfect,  WOULD + Verbo Base

"If I had studied medicine, I would be a doctor now."
→ (No estudié medicina en el pasado → ahora no soy médico)

PRESENTE IRREAL  →  consecuencia en el PASADO:
IF + Past Simple,  WOULD HAVE + Participio

"If I weren't so lazy, I would have finished it by now."
→ (Soy vago [presente] → no lo terminé [pasado])
\end{verbatim}

\begin{center}\rule{0.5\linewidth}{0.5pt}\end{center}

\section{⚠️ Errores Clásicos con
Condicionales}\label{errores-cluxe1sicos-con-condicionales}

\begin{verbatim}
❌  If I would have money, I would travel.
✅  If I had money, I would travel.

❌  If I would study, I will pass.
✅  If I study, I will pass.

❌  If she would have called...
✅  If she had called...

❌  I would have went if I had known.
✅  I would have gone if I had known.
\end{verbatim}

\begin{center}\rule{0.5\linewidth}{0.5pt}\end{center}

\section{📝 Mini Test Mental}\label{mini-test-mental}

\begin{quote}
Elige el condicional correcto: 1. \emph{``\_\_\_ you heat ice, it
melts.''} 2. \emph{``\_\_\_ I were rich, I would buy a yacht.''} 3.
\emph{``\_\_\_ it rains tomorrow, we'll cancel.''} 4. \emph{``\_\_\_ she
had studied, she would have passed.''}

\emph{(Respuestas: 1-If/When · 2-If · 3-If · 4-If)}
\end{quote}

\chapter{CAPÍTULO 15 - Relative
Clauses}\label{capuxedtulo-15---relative-clauses}

\emph{Las frases que dan información extra}

\begin{center}\rule{0.5\linewidth}{0.5pt}\end{center}

\begin{quote}
\subsection{🎯 Idea Central}\label{idea-central-13}

Las relative clauses añaden información sobre un sustantivo usando
pronombres relativos. La gran pregunta es: ¿esa información es
\textbf{ESENCIAL} o \textbf{EXTRA}?
\end{quote}

\begin{center}\rule{0.5\linewidth}{0.5pt}\end{center}

\section{📐 Los Pronombres Relativos}\label{los-pronombres-relativos}

\begin{verbatim}
WHO    →  personas
WHICH  →  cosas / animales
THAT   →  personas o cosas (uso informal/coloquial)
WHOSE  →  posesión (cuyo/a — de quien / de lo que)
WHERE  →  lugares
WHEN   →  tiempo
WHY    →  razón (solo con "the reason")
\end{verbatim}

\begin{center}\rule{0.5\linewidth}{0.5pt}\end{center}

\section{📐 Defining vs Non-Defining}\label{defining-vs-non-defining}

\begin{center}\rule{0.5\linewidth}{0.5pt}\end{center}

\subsubsection{🔵 DEFINING --- Información ESENCIAL (sin
comas)}\label{defining-informaciuxf3n-esencial-sin-comas}

\begin{verbatim}
"The woman WHO called you is my sister."
→  Sin esa información, no sabes de qué mujer hablas.
→  La cláusula define / identifica al sustantivo.
→  NO lleva comas.
→  Puedes usar THAT en lugar de WHO o WHICH.
→  Puedes OMITIR el pronombre si es objeto:
   "The book (that) I bought is great."  ✅
\end{verbatim}

\begin{center}\rule{0.5\linewidth}{0.5pt}\end{center}

\subsection{🔴 NON-DEFINING --- Información EXTRA (con
comas)}\label{non-defining-informaciuxf3n-extra-con-comas}

\begin{verbatim}
"My sister, WHO lives in London, called you."
→  Ya sabes de qué hermana hablas. La info es adicional.
→  Entre comas. Si se elimina, la frase sigue siendo correcta.
→  NO puedes usar THAT.
→  NO puedes omitir el pronombre relativo.
→  En lenguaje oral → pausa antes y después.
\end{verbatim}

\begin{center}\rule{0.5\linewidth}{0.5pt}\end{center}

\section{💡 Tabla Comparativa}\label{tabla-comparativa}

\begin{longtable}[]{@{}lll@{}}
\toprule\noalign{}
& Defining & Non-Defining \\
\midrule\noalign{}
\endhead
\bottomrule\noalign{}
\endlastfoot
Comas & ❌ Sin comas & ✅ Con comas \\
THAT & ✅ Se puede usar & ❌ No se puede \\
Omitir pronombre & ✅ Si es objeto & ❌ Nunca \\
Sin ella, la frase\ldots{} & pierde el sentido & sigue siendo
correcta \\
\end{longtable}

\begin{center}\rule{0.5\linewidth}{0.5pt}\end{center}

\section{✅ Ejemplos Comparados}\label{ejemplos-comparados-1}

\begin{verbatim}
DEFINING:
"The man who stole the car was arrested."
→  (qué hombre — el que robó el coche)

NON-DEFINING:
"John, who stole the car, was arrested."
→  (ya sabemos que es John, añadimos info extra)

DEFINING:
"The book that I recommended is a bestseller."

NON-DEFINING:
"'Don Quixote', which was written by Cervantes, is a classic."
\end{verbatim}

\begin{center}\rule{0.5\linewidth}{0.5pt}\end{center}

\section{⚠️ Errores Frecuentes}\label{errores-frecuentes-1}

\begin{verbatim}
❌  The book which I bought it is great.
✅  The book which I bought is great.
    (el relativo YA hace la función de objeto, no añadas "it")

❌  My brother, that lives in Paris, is a chef.
✅  My brother, who lives in Paris, is a chef.
    (non-defining → no uses THAT)

❌  The city where I visited was beautiful.
✅  The city which/that I visited was beautiful.
    (WHERE = en donde, aquí es objeto directo, no lugar)
    ✅  The city where I was born was beautiful.  (aquí sí es correcto)
\end{verbatim}

\chapter{CAPÍTULO 16. Wish \& If
Only}\label{capuxedtulo-16.-wish-if-only}

\emph{El club de los que lamentan}

\begin{center}\rule{0.5\linewidth}{0.5pt}\end{center}

\begin{quote}
\subsection{🎯 Idea Central}\label{idea-central-14}

\textbf{WISH} expresa deseos sobre situaciones \textbf{contrarias a la
realidad.} Es la gramática del ``ojalá'' y el ``si tan solo\ldots{}''
\end{quote}

\begin{center}\rule{0.5\linewidth}{0.5pt}\end{center}

\section{📐 Las Tres Estructuras de
WISH}\label{las-tres-estructuras-de-wish}

\begin{verbatim}
WISH  +  Past Simple
→  Deseo sobre el PRESENTE (algo irreal ahora)
"I wish I spoke Spanish."     →  Ojalá hablara español. (no hablo)
"I wish I had more time."     →  Ojalá tuviera más tiempo. (no tengo)
"I wish it were Friday."      →  Ojalá fuera viernes. (no lo es)

WISH  +  Past Perfect
→  Deseo / lamento sobre el PASADO (ya no se puede cambiar)
"I wish I had studied more."  →  Ojalá hubiera estudiado más. (no estudié)
"I wish I hadn't said that."  →  Ojalá no lo hubiera dicho.

WISH  +  WOULD
→  Deseo de cambio de comportamiento (queja / irritación)
"I wish you would stop smoking."  →  Ojalá dejaras de fumar.
"I wish it would stop raining."   →  Ojalá dejara de llover.
⚠️  No uses WISH + WOULD para ti mismo:
❌  I wish I would study more.  →  ✅  I wish I studied more.
\end{verbatim}

\begin{center}\rule{0.5\linewidth}{0.5pt}\end{center}

\section{💡 Truco --- WISH vs HOPE}\label{truco-wish-vs-hope}

\begin{verbatim}
WISH  →  la realidad es diferente a lo que quieres
         (es irreal, imposible o muy improbable ahora)
"I wish I were taller."   →  (no lo soy, es un hecho)

HOPE  →  la situación es posible, esperas que ocurra
"I hope it doesn't rain." →  (es posible que no llueva)
"I hope you feel better." →  (es posible que te mejores)
\end{verbatim}

\begin{center}\rule{0.5\linewidth}{0.5pt}\end{center}

\section{🔥 IF ONLY --- El WISH más
dramático}\label{if-only-el-wish-muxe1s-dramuxe1tico}

\begin{verbatim}
IF ONLY tiene las mismas estructuras que WISH
pero con mucho más énfasis emocional:

"If only I had more money!"     →  ¡Si tan solo tuviera más dinero!
"If only I hadn't trusted him!" →  ¡Si tan solo no le hubiera creído!
"If only you would listen!"     →  ¡Si tan solo escucharas!
\end{verbatim}

\chapter{CAPÍTULO 17 Phrasal Verbs}\label{capuxedtulo-17-phrasal-verbs}

\emph{El código secreto del inglés real}

\begin{center}\rule{0.5\linewidth}{0.5pt}\end{center}

\begin{quote}
\subsection{🎯 Idea Central}\label{idea-central-15}

Verbo + partícula = significado completamente nuevo. \textbf{No los
traduzcas literalmente. Agrúpalos por partículas.} Aprende el
significado de la partícula y entenderás grupos enteros a la vez.
\end{quote}

\begin{center}\rule{0.5\linewidth}{0.5pt}\end{center}

\section{💡 El Significado de las
Partículas}\label{el-significado-de-las-partuxedculas}

\begin{center}\rule{0.5\linewidth}{0.5pt}\end{center}

\subsection{⬆️ UP --- Completar, aumentar, mejorar,
finalizar}\label{up-completar-aumentar-mejorar-finalizar}

\begin{verbatim}
give UP      →  rendirse           "Don't give up!"
wake UP      →  despertar          "Wake up! It's late."
grow UP      →  crecer             "I grew up in Madrid."
end UP       →  terminar siendo    "I ended up staying home."
set UP       →  establecer/crear   "She set up her own company."
come UP with →  ocurrírsele algo   "He came up with a great idea."
build UP     →  desarrollar        "She built up her confidence."
make UP      →  inventar / maquillarse "He made up an excuse."
turn UP      →  aparecer / subir volumen "She turned up late."
clean UP     →  limpiar del todo   "Clean up your room!"
\end{verbatim}

\begin{center}\rule{0.5\linewidth}{0.5pt}\end{center}

\subsection{⬇️ OUT --- Salir, descubrir, completar,
extinguir}\label{out-salir-descubrir-completar-extinguir}

\begin{verbatim}
find OUT     →  descubrir          "I found out the truth."
run OUT of   →  quedarse sin       "We've run out of milk."
work OUT     →  hacer ejercicio / resolver "It worked out fine."
turn OUT     →  resultar           "It turned out to be great."
carry OUT    →  llevar a cabo      "They carried out the plan."
point OUT    →  señalar            "She pointed out the mistake."
figure OUT   →  entender/resolver  "I can't figure it out."
put OUT      →  apagar / extinguir "Put out the fire!"
stand OUT    →  destacar           "She always stands out."
\end{verbatim}

\begin{center}\rule{0.5\linewidth}{0.5pt}\end{center}

\subsection{❌ OFF --- Separación, cancelación, salida,
inicio}\label{off-separaciuxf3n-cancelaciuxf3n-salida-inicio}

\begin{verbatim}
turn OFF     →  apagar             "Turn off the light."
put OFF      →  posponer           "Don't put it off!"
take OFF     →  despegar / quitarse "The plane took off."
call OFF     →  cancelar           "They called off the meeting."
set OFF      →  partir / salir     "We set off at 7am."
pay OFF      →  merecer la pena    "Hard work always pays off."
show OFF     →  presumir           "Stop showing off!"
\end{verbatim}

\begin{center}\rule{0.5\linewidth}{0.5pt}\end{center}

\subsection{▶️ ON --- Continuación, conexión,
avance}\label{on-continuaciuxf3n-conexiuxf3n-avance}

\begin{verbatim}
carry ON     →  continuar          "Carry on with your work."
get on with  →  llevarse bien / continuar "I get on with my boss."
hold ON      →  esperar            "Hold on a moment."
move ON      →  seguir adelante    "It's time to move on."
rely ON      →  depender de        "You can rely on me."
take ON      →  asumir / contratar "She took on a new project."
\end{verbatim}

\begin{center}\rule{0.5\linewidth}{0.5pt}\end{center}

\subsection{🔄 INTO --- Transformación, investigación,
choque}\label{into-transformaciuxf3n-investigaciuxf3n-choque}

\begin{verbatim}
turn INTO    →  convertirse en     "Water turns into ice."
break INTO   →  entrar a la fuerza "They broke into the house."
look INTO    →  investigar         "Police are looking into it."
run INTO     →  encontrarse con    "I ran into my old teacher."
\end{verbatim}

\begin{center}\rule{0.5\linewidth}{0.5pt}\end{center}

\subsection{🔙 BACK --- Retorno, devolución,
reciprocidad}\label{back-retorno-devoluciuxf3n-reciprocidad}

\begin{verbatim}
get BACK     →  volver / recuperar "I got back home late."
give BACK    →  devolver           "Give back what you took."
call BACK    →  devolver la llamada "I'll call you back."
look BACK    →  mirar atrás / recordar "Looking back, I understand."
pay BACK     →  devolver dinero    "I'll pay you back tomorrow."
\end{verbatim}

\begin{center}\rule{0.5\linewidth}{0.5pt}\end{center}

\section{⚠️ Phrasal Verbs Separables vs
Inseparables}\label{phrasal-verbs-separables-vs-inseparables}

\begin{verbatim}
SEPARABLES  →  puedes poner el objeto entre el verbo y la partícula:
"Turn off the light."    ✅
"Turn the light off."    ✅
"Turn it off."           ✅  (pronombre: SIEMPRE en el medio)
"Turn off it."           ❌  (pronombre nunca al final)

INSEPARABLES  →  el objeto va siempre después de la partícula:
"Look after the children."  ✅
"Look the children after."  ❌
"Look after them."          ✅
\end{verbatim}

\chapter{CAPÍTULO 18 Linking Words \&
Connectors}\label{capuxedtulo-18-linking-words-connectors}

\emph{El pegamento de tus ideas --- Esencial para el B2}

\begin{quote}
\subsection{🎯 Idea Central}\label{idea-central-16}

Los conectores convierten una lista de ideas en un texto que
\textbf{FLUYE naturalmente.} Sin ellos, tienes frases. Con ellos, tienes
un texto.
\end{quote}

\begin{center}\rule{0.5\linewidth}{0.5pt}\end{center}

\section{📐 Categorías de Conectores}\label{categoruxedas-de-conectores}

\begin{center}\rule{0.5\linewidth}{0.5pt}\end{center}

\subsection{➕ AÑADIR IDEAS}\label{auxf1adir-ideas}

\begin{verbatim}
Furthermore      →  Además (más formal que "also")
Moreover         →  Es más / Además (aún más formal)
In addition      →  Adicionalmente
What is more     →  Es más
Not only... but also → No solo... sino también
Also             →  También (menos formal)
Besides          →  Además / Aparte de eso
\end{verbatim}

\begin{center}\rule{0.5\linewidth}{0.5pt}\end{center}

\subsection{↔️ CONTRASTAR Y CONCEDER}\label{contrastar-y-conceder}

\begin{verbatim}
However          →  Sin embargo
Nevertheless     →  No obstante (más formal)
Nonetheless      →  No obstante (muy formal)
On the other hand→  Por otro lado
In contrast      →  En contraste
Although         →  Aunque (+ frase completa)
Though           →  Aunque (más informal, puede ir al final)
Even though      →  Aunque (con más énfasis)
Despite          →  A pesar de (+ nombre o -ING)
In spite of      →  A pesar de (+ nombre o -ING)
Whereas          →  Mientras que / En cambio
While            →  Mientras que (contraste)
Yet              →  Sin embargo (formal, después de coma)
\end{verbatim}

\begin{quote}
💡 \textbf{ALTHOUGH vs DESPITE --- El contraste más evaluado:}

\begin{verbatim}
ALTHOUGH  +  frase completa (sujeto + verbo):
"Although it was raining, we went out."

DESPITE / IN SPITE OF  +  sustantivo o -ING:
"Despite the rain, we went out."
"In spite of raining heavily, we went out."
\end{verbatim}
\end{quote}

\begin{center}\rule{0.5\linewidth}{0.5pt}\end{center}

\subsection{🎯 CAUSA Y EFECTO}\label{causa-y-efecto}

\begin{verbatim}
Therefore        →  Por lo tanto (consecuencia lógica)
As a result      →  Como resultado
Consequently     →  Consecuentemente
Hence            →  De ahí que (muy formal)
Thus             →  Así pues (formal)
Due to           →  Debido a (+ nombre)
Owing to         →  Debido a (formal, + nombre)
Because of       →  A causa de (+ nombre)
Since / As       →  Ya que / Dado que (+ frase)
So               →  Así que (informal)
\end{verbatim}

\begin{center}\rule{0.5\linewidth}{0.5pt}\end{center}

\subsection{📝 EJEMPLIFICAR}\label{ejemplificar}

\begin{verbatim}
For example      →  Por ejemplo
For instance     →  Por ejemplo (más formal)
Such as          →  Como (por ejemplo): "fruits such as apples"
Including        →  Incluyendo
Namely           →  A saber / Es decir (lista específica)
In particular    →  En particular
Specifically     →  Específicamente
\end{verbatim}

\begin{center}\rule{0.5\linewidth}{0.5pt}\end{center}

\subsection{📋 SECUENCIAR IDEAS}\label{secuenciar-ideas}

\begin{verbatim}
First of all     →  En primer lugar
To begin with    →  Para empezar
Firstly / Secondly / Finally → Primero / Segundo / Finalmente
Then / Next      →  Luego / A continuación
Subsequently     →  Posteriormente (formal)
Meanwhile        →  Mientras tanto
Eventually       →  Finalmente / al cabo de un tiempo
\end{verbatim}

\begin{center}\rule{0.5\linewidth}{0.5pt}\end{center}

\subsection{🏁 CONCLUIR}\label{concluir}

\begin{verbatim}
To conclude      →  Para concluir
To sum up        →  Para resumir
In conclusion    →  En conclusión
All in all       →  En definitiva
On the whole     →  En general
Overall          →  En general / Globalmente
Taking everything into consideration → Considerándolo todo
\end{verbatim}

\begin{center}\rule{0.5\linewidth}{0.5pt}\end{center}

\section{💡 Tabla Rápida --- Nivel de
Formalidad}\label{tabla-ruxe1pida-nivel-de-formalidad}

\begin{longtable}[]{@{}lll@{}}
\toprule\noalign{}
Conector & Equivalente informal & Nivel \\
\midrule\noalign{}
\endhead
\bottomrule\noalign{}
\endlastfoot
Furthermore & Also & ⭐⭐⭐ Muy formal \\
Nevertheless & But & ⭐⭐⭐ Muy formal \\
Therefore & So & ⭐⭐⭐ Formal \\
However & But & ⭐⭐ Formal \\
Although & But & ⭐⭐ Semi-formal \\
In addition & And & ⭐⭐ Semi-formal \\
\end{longtable}

\chapter{CAPÍTULO 19 B2 Writing --- Guía
Completa}\label{capuxedtulo-19-b2-writing-guuxeda-completa}

\emph{Todo lo que necesitas para brillar en el writing}

\begin{center}\rule{0.5\linewidth}{0.5pt}\end{center}

\begin{quote}
\subsection{🎯 Idea Central}\label{idea-central-17}

El writing del B2 evalúa: \textbf{organización, variedad gramatical,
vocabulario avanzado, registro adecuado y cohesión.} No es solo escribir
sin errores --- es escribir con \textbf{impacto.}
\end{quote}

\begin{center}\rule{0.5\linewidth}{0.5pt}\end{center}

\section{📧 FORMAL EMAIL / LETTER}\label{formal-email-letter}

\begin{verbatim}
ESTRUCTURA:
┌─────────────────────────────────────────────┐
│  Saludo formal                               │
│  Párrafo 1: motivo de la carta/email         │
│  Párrafo 2: desarrollo / detalles            │
│  Párrafo 3: petición / solución              │
│  Despedida formal                            │
└─────────────────────────────────────────────┘

SALUDOS:
Dear Mr / Mrs / Ms [Apellido],   →  (si conoces el nombre)
Dear Sir / Madam,                →  (si no conoces el nombre)

PRIMERA FRASE:
"I am writing with regard to..."
"I am writing to enquire about..."
"I am writing to complain about..."
"I am writing in response to your advertisement..."
"Further to our recent conversation, I am writing to..."

FRASES ÚTILES EN EL CUERPO:
"I would be grateful if you could..."
"I would appreciate it if..."
"Please do not hesitate to contact me..."
"I look forward to hearing from you."
"I would like to draw your attention to..."
"I am afraid I must express my dissatisfaction with..."

DESPEDIDAS:
Yours sincerely,   →  (si usaste nombre: Dear Mr Smith)
Yours faithfully,  →  (si usaste Dear Sir/Madam)
\end{verbatim}

\begin{center}\rule{0.5\linewidth}{0.5pt}\end{center}

\section{📰 ESSAY --- Opinion /
Argumentative}\label{essay-opinion-argumentative}

\begin{verbatim}
ESTRUCTURA:
┌─────────────────────────────────────────────┐
│  Introducción: presenta el tema              │
│               enuncia tu postura             │
│  Párrafo 2: primer argumento + ejemplo       │
│  Párrafo 3: segundo argumento + ejemplo      │
│  Párrafo 4: argumento contrario + refutación │
│  Conclusión: resume y reafirma postura       │
└─────────────────────────────────────────────┘

INTRODUCCIÓN:
"It is widely believed that..."
"In recent years, the issue of... has become increasingly relevant."
"There is no doubt that... is one of the most controversial topics today."
"The question of whether... has long been debated."

DESARROLLAR ARGUMENTOS:
"One of the main advantages of... is..."
"A key argument in favour of... is..."
"Furthermore, it should be noted that..."
"Research has shown that..."
"This is clearly demonstrated by..."

ARGUMENTO CONTRARIO:
"Critics argue that... However, this view fails to consider..."
"While it is true that..., the fact remains that..."
"Opponents claim that... Nevertheless..."

CONCLUSIÓN:
"Taking everything into consideration, I firmly believe that..."
"All things considered, it seems clear that..."
"On balance, the advantages clearly outweigh the disadvantages."
\end{verbatim}

\begin{center}\rule{0.5\linewidth}{0.5pt}\end{center}

\section{📖 ARTICLE}\label{article}

\begin{verbatim}
ESTRUCTURA:
┌─────────────────────────────────────────────┐
│  Título llamativo (pregunta o gancho)        │
│  Introducción que engancha al lector         │
│  Párrafos de desarrollo                      │
│  Conclusión con llamada a la acción          │
└─────────────────────────────────────────────┘

CARACTERÍSTICAS:
- Tono puede ser semi-formal o informal
- Dirígete al lector directamente:
  "Have you ever wondered...?"
  "You might be surprised to learn that..."
  "Think about the last time you..."
- Usa anécdotas y ejemplos concretos
- Título con juego de palabras o pregunta
\end{verbatim}

\begin{center}\rule{0.5\linewidth}{0.5pt}\end{center}

\section{📝 REPORT}\label{report}

\begin{verbatim}
ESTRUCTURA:
┌─────────────────────────────────────────────┐
│  TITLE: Report on [tema]                    │
│  INTRODUCTION: objetivo del informe          │
│  FINDINGS: resultados / hallazgos            │
│  CONCLUSIONS: conclusiones                   │
│  RECOMMENDATIONS: recomendaciones            │
└─────────────────────────────────────────────┘

FRASES TÍPICAS:
"The aim of this report is to..."
"The findings suggest that..."
"It is recommended that..."
"Based on the data collected, it can be concluded that..."
"The evidence indicates that..."
\end{verbatim}

\begin{center}\rule{0.5\linewidth}{0.5pt}\end{center}

\section{🔥 Los 5 Mandamientos del B2
Writing}\label{los-5-mandamientos-del-b2-writing}

\begin{verbatim}
1.  Varía tu vocabulario — nunca repitas la misma palabra dos veces
2.  Usa al menos 3 tiempos verbales diferentes en tu texto
3.  Incluye estructuras pasivas y modales
4.  Conecta cada párrafo con un conector apropiado
5.  Adapta el REGISTRO al tipo de texto (formal/informal)
\end{verbatim}

\chapter{CAPÍTULO 20 Vocabulario Avanzado
B2}\label{capuxedtulo-20-vocabulario-avanzado-b2}

\emph{Las palabras que el examen espera ver}

\begin{center}\rule{0.5\linewidth}{0.5pt}\end{center}

\begin{quote}
\section{Idea Central}\label{idea-central-18}

El B2 penaliza la repetición. La riqueza léxica es uno de los criterios
de evaluación. \textbf{Aprende familias de sinónimos, no palabras
aisladas.}
\end{quote}

\begin{center}\rule{0.5\linewidth}{0.5pt}\end{center}

\section{Sinónimos Esenciales --- Los Más
Evaluados}\label{sinuxf3nimos-esenciales-los-muxe1s-evaluados}

\begin{center}\rule{0.5\linewidth}{0.5pt}\end{center}

\subsection{En lugar de GOOD:}\label{en-lugar-de-good}

\begin{verbatim}
excellent, outstanding, remarkable, beneficial,
advantageous, effective, valuable, significant,
commendable, impressive, worthwhile, productive
\end{verbatim}

\subsection{En lugar de BAD:}\label{en-lugar-de-bad}

\begin{verbatim}
detrimental, harmful, damaging, negative,
problematic, concerning, disadvantageous,
counterproductive, adverse, unfavourable
\end{verbatim}

\subsection{En lugar de BIG /
IMPORTANT:}\label{en-lugar-de-big-important}

\begin{verbatim}
significant, substantial, considerable, extensive,
enormous, vast, major, large-scale, crucial, vital,
essential, fundamental, key, critical, paramount
\end{verbatim}

\subsection{En lugar de THINK:}\label{en-lugar-de-think}

\begin{verbatim}
believe, consider, argue, suggest, claim,
maintain, contend, feel, reckon (informal),
assert, acknowledge, recognize, propose
\end{verbatim}

\subsection{En lugar de SHOW:}\label{en-lugar-de-show}

\begin{verbatim}
demonstrate, reveal, indicate, suggest,
highlight, illustrate, prove, confirm,
reflect, imply, signify, establish
\end{verbatim}

\subsection{En lugar de INCREASE:}\label{en-lugar-de-increase}

\begin{verbatim}
rise, grow, expand, surge, soar, escalate,
climb, develop, improve, enhance, boost
\end{verbatim}

\subsection{En lugar de DECREASE:}\label{en-lugar-de-decrease}

\begin{verbatim}
fall, drop, decline, reduce, diminish,
shrink, plummet, deteriorate, worsen
\end{verbatim}

\begin{center}\rule{0.5\linewidth}{0.5pt}\end{center}

\section{Vocabulario Temático --- Los Temas del
B2}\label{vocabulario-temuxe1tico-los-temas-del-b2}

\begin{center}\rule{0.5\linewidth}{0.5pt}\end{center}

\subsection{MEDIO AMBIENTE}\label{medio-ambiente}

\begin{verbatim}
climate change          →  cambio climático
global warming          →  calentamiento global
greenhouse gases        →  gases de efecto invernadero
renewable energy        →  energía renovable
carbon footprint        →  huella de carbono
biodiversity            →  biodiversidad
deforestation           →  deforestación
sustainable             →  sostenible
environmentally friendly→  respetuoso con el medioambiente
\end{verbatim}

\subsection{TECNOLOGÍA}\label{tecnologuxeda}

\begin{verbatim}
artificial intelligence →  inteligencia artificial
data privacy            →  privacidad de datos
social media            →  redes sociales
digital divide          →  brecha digital
cybersecurity           →  ciberseguridad
remote work             →  trabajo a distancia
automation              →  automatización
\end{verbatim}

\subsection{SALUD Y SOCIEDAD}\label{salud-y-sociedad}

\begin{verbatim}
mental health           →  salud mental
well-being              →  bienestar
healthcare system       →  sistema sanitario
work-life balance       →  equilibrio vida-trabajo
inequality              →  desigualdad
social cohesion         →  cohesión social
\end{verbatim}

\chapter{CAPÍTULO 21 Falsos Amigos}\label{capuxedtulo-21-falsos-amigos}

\emph{Las palabras que te traicionan}

\begin{center}\rule{0.5\linewidth}{0.5pt}\end{center}

\begin{quote}
\section{Idea Central}\label{idea-central-19}

Los falsos amigos son palabras que \textbf{se parecen} en español e
inglés pero significan cosas \textbf{completamente diferentes.}
Apréndetelos: pueden cambiar el sentido de toda una frase.
\end{quote}

\begin{center}\rule{0.5\linewidth}{0.5pt}\end{center}

\section{Los Falsos Amigos Más
Peligrosos}\label{los-falsos-amigos-muxe1s-peligrosos}

\begin{longtable}[]{@{}
  >{\raggedright\arraybackslash}p{(\linewidth - 6\tabcolsep) * \real{0.1882}}
  >{\raggedright\arraybackslash}p{(\linewidth - 6\tabcolsep) * \real{0.2118}}
  >{\raggedright\arraybackslash}p{(\linewidth - 6\tabcolsep) * \real{0.2471}}
  >{\raggedright\arraybackslash}p{(\linewidth - 6\tabcolsep) * \real{0.3529}}@{}}
\toprule\noalign{}
\begin{minipage}[b]{\linewidth}\raggedright
Palabra inglesa
\end{minipage} & \begin{minipage}[b]{\linewidth}\raggedright
Parece significar
\end{minipage} & \begin{minipage}[b]{\linewidth}\raggedright
Realmente significa
\end{minipage} & \begin{minipage}[b]{\linewidth}\raggedright
Cómo se dice lo que pensabas
\end{minipage} \\
\midrule\noalign{}
\endhead
\bottomrule\noalign{}
\endlastfoot
\textbf{actually} & actualmente & en realidad & currently, nowadays \\
\textbf{eventually} & eventualmente & finalmente, al final & possibly,
perhaps \\
\textbf{embarrassed} & embarazada & avergonzado/a & pregnant \\
\textbf{sensible} & sensible (emocional) & sensato, razonable &
sensitive \\
\textbf{sensitive} & sensitivo & sensible (emocional) & --- \\
\textbf{assist} & asistir (ir) & ayudar & attend \\
\textbf{attend} & atender & asistir (ir a) & serve, help \\
\textbf{library} & librería & biblioteca & bookshop \\
\textbf{fabric} & fábrica & tela, tejido & factory \\
\textbf{carpet} & carpeta & moqueta, alfombra & folder \\
\textbf{success} & suceso & éxito & event, incident \\
\textbf{parents} & parientes & padres & relatives \\
\textbf{college} & colegio (infantil) & universidad/facultad & school \\
\textbf{pretend} & pretender & fingir, simular & claim, intend \\
\textbf{realize} & realizar & darse cuenta & carry out, achieve \\
\textbf{exit} & éxito & salida & success \\
\textbf{sympathetic} & simpático & comprensivo, empático & nice,
friendly \\
\textbf{comprehensive} & comprensivo & exhaustivo, completo &
understanding \\
\textbf{lecture} & lectura & clase magistral, conferencia & reading \\
\textbf{record} & recordar & grabar / récord & remember \\
\textbf{introduce} & introducir & presentar a alguien & insert, put
in \\
\textbf{large} & largo & grande & long \\
\textbf{actually} & actualmente & de hecho, en realidad & currently \\
\end{longtable}

\begin{center}\rule{0.5\linewidth}{0.5pt}\end{center}

\section{💡 Los Tres Más Usados --- Memorízalos
Ya}\label{los-tres-muxe1s-usados-memoruxedzalos-ya}

\begin{verbatim}
ACTUALLY vs CURRENTLY:
"Actually, I don't agree."  →  En realidad, no estoy de acuerdo.
"Currently, I work in IT."  →  Actualmente, trabajo en informática.

EMBARRASSED vs PREGNANT:
"I was so embarrassed."     →  Estaba muy avergonzado/a.
"She is pregnant."          →  Está embarazada.

SENSIBLE vs SENSITIVE:
"She's very sensible."      →  Es muy sensata / razonable.
"He's very sensitive."      →  Es muy sensible (se emociona fácilmente).
\end{verbatim}

\chapter{CAPÍTULO 22. Pronunciation \& Speaking
B2}\label{capuxedtulo-22.-pronunciation-speaking-b2}

\emph{Sonar bien sin nacer en Londres}

\begin{center}\rule{0.5\linewidth}{0.5pt}\end{center}

\begin{quote}
\subsection{🎯 Idea Central}\label{idea-central-20}

No necesitas acento perfecto. Necesitas ser \textbf{comprensible y
fluido.} El B2 evalúa fluidez, coherencia y rango de vocabulario --- no
el acento.
\end{quote}

\begin{center}\rule{0.5\linewidth}{0.5pt}\end{center}

\section{Los Sonidos que No Existen en
Español}\label{los-sonidos-que-no-existen-en-espauxf1ol}

\begin{center}\rule{0.5\linewidth}{0.5pt}\end{center}

\subsection{El Sonido /θ/ --- TH sorda (think, three,
through)}\label{el-sonido-ux3b8-th-sorda-think-three-through}

\begin{verbatim}
🧠 Pon la punta de la lengua ENTRE los dientes y sopla suavemente.
   Piensa que estás "empezando a cecear".

think   →  no es "tink" ni "sink"
three   →  no es "tri" ni "sri"
through →  no es "tru" ni "sru"
\end{verbatim}

\subsection{El Sonido /ð/ --- TH sonora (the, this, that,
brother)}\label{el-sonido-uxf0-th-sonora-the-this-that-brother}

\begin{verbatim}
🧠 Igual que el anterior pero CON VOZ (vibración en la garganta).

the     →  no es "de" ni "ze"
this    →  no es "dis"
brother →  no es "brozer"
\end{verbatim}

\subsection{El Sonido /æ/ --- A de cat, man,
black}\label{el-sonido-uxe6-a-de-cat-man-black}

\begin{verbatim}
🧠 Entre la A y la E españolas, con la boca más abierta.
   Como si empezaras a decir "e" pero con la boca de "a".

cat   →  /kæt/  (no como la "a" española)
man   →  /mæn/
black →  /blæk/
\end{verbatim}

\begin{center}\rule{0.5\linewidth}{0.5pt}\end{center}

\section{El Caos del -OUGH --- 6 Pronunciaciones
Diferentes}\label{el-caos-del--ough-6-pronunciaciones-diferentes}

\begin{verbatim}
thought  →  /θɔːt/   (aunque, pensó)
though   →  /ðəʊ/    (aunque)
through  →  /θruː/   (a través)
thorough →  /ˈθʌrə/  (exhaustivo)
tough    →  /tʌf/    (duro)
cough    →  /kɒf/    (tos)

🧠 No hay regla. Acéptalo. Memorízalos uno a uno.
\end{verbatim}

\begin{center}\rule{0.5\linewidth}{0.5pt}\end{center}

\section{Word Stress --- El Acento que Cambia el
Significado}\label{word-stress-el-acento-que-cambia-el-significado}

\begin{verbatim}
Algunas palabras cambian significado según dónde va el acento:

REcord  (nombre) → un disco, un récord
reCORD  (verbo)  → grabar

PREsent (nombre/adj) → regalo / presente
preSENT (verbo)      → presentar

PERmit  (nombre) → permiso (documento)
perMIT  (verbo)  → permitir

OBject  (nombre) → objeto
obJECT  (verbo)  → objetar
\end{verbatim}

\begin{center}\rule{0.5\linewidth}{0.5pt}\end{center}

\section{B2 Speaking --- Lo Que Dices Cuando No Sabes Qué
Decir}\label{b2-speaking-lo-que-dices-cuando-no-sabes-quuxe9-decir}

\begin{center}\rule{0.5\linewidth}{0.5pt}\end{center}

\subsection{Ganar Tiempo --- Fillers}\label{ganar-tiempo-fillers}

\begin{verbatim}
"That's a good question..."
"Let me think about that for a moment..."
"Well, to be honest..."
"I'd say that..."
"As far as I know..."
"If I'm not mistaken..."
"It's difficult to say, but..."
"I've never really thought about it, but I'd say..."
\end{verbatim}

\begin{center}\rule{0.5\linewidth}{0.5pt}\end{center}

\subsection{Dar Opiniones}\label{dar-opiniones}

\begin{verbatim}
"In my opinion / view..."
"From my perspective..."
"As far as I'm concerned..."
"I strongly believe that..."
"I tend to think that..."
"Personally, I feel that..."
"It seems to me that..."
"I'm convinced that..."
\end{verbatim}

\begin{center}\rule{0.5\linewidth}{0.5pt}\end{center}

\subsection{Especular --- Muy Importante en
B2}\label{especular-muy-importante-en-b2}

\begin{verbatim}
"It looks as if / though..."
"It seems to me that..."
"I imagine that..."
"I would guess that..."
"She might be... / She could be..."
"Perhaps / Possibly..."
"I wouldn't be surprised if..."
"There's a chance that..."
\end{verbatim}

\begin{center}\rule{0.5\linewidth}{0.5pt}\end{center}

\subsection{Reformular --- Cuando No Recuerdas la
Palabra}\label{reformular-cuando-no-recuerdas-la-palabra}

\begin{verbatim}
"What I mean is..."
"In other words..."
"What I'm trying to say is..."
"It's a bit like..."
"It's similar to..."
"It's the kind of thing that..."
"It's used for... -ing"     ←  para describir un objeto
\end{verbatim}

\begin{center}\rule{0.5\linewidth}{0.5pt}\end{center}

\subsection{Estructurar Tu Respuesta}\label{estructurar-tu-respuesta}

\begin{verbatim}
INTRODUCIR:   "First of all... / To begin with..."
AÑADIR:       "Furthermore... / What's more... / On top of that..."
CONTRASTAR:   "On the other hand... / Having said that..."
EJEMPLIFICAR: "For instance... / A good example of this is..."
CONCLUIR:     "All things considered... / On balance... / Overall..."
\end{verbatim}

\begin{center}\rule{0.5\linewidth}{0.5pt}\end{center}

\section{Los 5 Errores Más Comunes en el Speaking
B2}\label{los-5-errores-muxe1s-comunes-en-el-speaking-b2}

\begin{verbatim}
1.  ❌  Traducir directamente del español → frases raras
    ✅  Piensa en ideas, no en traducciones palabra por palabra

2.  ❌  Respuestas de una sola frase
    ✅  Desarrolla, justifica, da ejemplos: mínimo 3-4 frases

3.  ❌  Usar solo "I think" para todo
    ✅  Varía: I believe / I feel / I'd argue / It seems to me

4.  ❌  Silencio cuando no sabes qué decir
    ✅  Usa fillers: "That's an interesting question..."

5.  ❌  Hablar demasiado rápido por nervios
    ✅  Despacio y claro > rápido e incomprensible
\end{verbatim}

\bookmarksetup{startatroot}

\chapter{RESUMEN FINAL --- El Mapa
Completo}\label{resumen-final-el-mapa-completo}

\section{MAPA GRAMATICAL}\label{mapa-gramatical}

\begin{verbatim}
╔══════════════════════════════════════════════════════════════╗
║                    MAPA GRAMATICAL COMPLETO                 ║
╠══════════════════════════════════════════════════════════════╣
║  📘 A1   To be · Pronombres · Present Simple                ║
║          Present Continuous · Stative Verbs                 ║
╠══════════════════════════════════════════════════════════════╣
║  📗 A2   Past Simple · Past Continuous                      ║
║          Future: Will / Going to / Present Continuous       ║
╠══════════════════════════════════════════════════════════════╣
║  📙 B1   Present Perfect · Present Perfect Continuous       ║
║          Past Perfect · Modal Verbs completos               ║
║          Passive Voice · Reported Speech                    ║
╠══════════════════════════════════════════════════════════════╣
║  📕 B2   Conditionals (0,1,2,3 + Mixed)                     ║
║          Relative Clauses · Wish & If Only                  ║
║          Modal Perfects · Phrasal Verbs                     ║
║          Linking Words · Writing Skills                     ║
║          Vocabulario avanzado · Speaking strategies         ║
╚══════════════════════════════════════════════════════════════╝
\end{verbatim}

\begin{center}\rule{0.5\linewidth}{0.5pt}\end{center}

\section{Plan de Estudio Recomendado}\label{plan-de-estudio-recomendado}

\begin{longtable}[]{@{}lll@{}}
\toprule\noalign{}
Semana & Foco & Capítulos \\
\midrule\noalign{}
\endhead
\bottomrule\noalign{}
\endlastfoot
1 & Bases y presente & 1, 2, 3, 4 \\
2 & Pasado y futuro & 5, 6, 7 \\
3 & Present Perfect & 8, 9, 10 \\
4 & Modales y Pasiva & 11, 12 \\
5 & Reported Speech & 13 \\
6 & Condicionales & 14 \\
7 & Relative, Wish & 15, 16 \\
8 & Léxico y Writing & 17, 18, 19, 20 \\
9 & Falsos amigos y Speaking & 21, 22 \\
10 & Repaso general + simulacro & Todos \\
\end{longtable}

\begin{center}\rule{0.5\linewidth}{0.5pt}\end{center}

\section{Chuleta de Emergencia --- Los 10 Puntos Más Evaluados en
B2}\label{chuleta-de-emergencia-los-10-puntos-muxe1s-evaluados-en-b2}

\begin{verbatim}
1.  Present Perfect vs Past Simple        → Capítulo 8
2.  Condicionales (especialmente 2 y 3)   → Capítulo 14
3.  Reported Speech con backshift         → Capítulo 13
4.  Passive Voice en todos los tiempos    → Capítulo 12
5.  Modal Perfects (should have, etc.)    → Capítulo 11
6.  Relative Clauses con/sin comas        → Capítulo 15
7.  Wish + tiempos correctos              → Capítulo 16
8.  Linking Words en writing              → Capítulo 18
9.  Registro formal en emails/essays      → Capítulo 19
10. Vocabulario variado (no repetir)      → Capítulo 20
\end{verbatim}

\begin{center}\rule{0.5\linewidth}{0.5pt}\end{center}

\begin{quote}
\section{💬 Nota Final del Autor}\label{nota-final-del-autor}

\emph{Este libro no te da el inglés. Te da las llaves.} \emph{Tú tienes
que abrir las puertas.}

\emph{Cada vez que cometas un error, recuerda:} \emph{es tu cerebro
recalibrando.} \emph{Los errores no son fracasos ---} \emph{son el
sonido que hace tu mente mientras crece.}

\emph{Estudia este libro, pero sobre todo:} \textbf{\emph{habla,
escribe, escucha, lee.}}

\emph{El B2 no es un nivel. Es una actitud.}

\subsection{Good luck. You've got this.
🚀}\label{good-luck.-youve-got-this.}
\end{quote}

\begin{center}\rule{0.5\linewidth}{0.5pt}\end{center}

\begin{verbatim}
📓 FIN DEL LIBRO DE NOTAS DEFINITIVO — A1 a B2
   Versión 2.0 — Revisada y Ampliada
   Hecho con 💙 por un hispanohablante
   que aprendió a las malas para que tú aprendas a las buenas.
\end{verbatim}

\bookmarksetup{startatroot}

\chapter*{References}\label{references}
\addcontentsline{toc}{chapter}{References}

\markboth{References}{References}

\phantomsection\label{refs}


\backmatter


\end{document}
